% SENTINEL-E: Anytime-Valid Sequential Detection for Streaming Illegal Dumping
% Surveillance.  Prepared for the Neural Computing and Applications topical
% collection "Emerging Trends in Smart Waste Monitoring and Environmental
% Surveillance".  For submission, replace the preamble with Springer's
% sn-jnl.cls; the body is class-independent.
\documentclass[11pt,a4paper]{article}

\usepackage[T1]{fontenc}
\usepackage[utf8]{inputenc}
\usepackage{lmodern}
\usepackage[margin=2.5cm]{geometry}
\usepackage{amsmath,amssymb,amsthm}
\usepackage{graphicx}
\usepackage{booktabs}
\usepackage{pifont}
\usepackage{caption}
\usepackage{subcaption}
\usepackage{microtype}
\usepackage[section]{placeins}
\usepackage{algorithm}
\usepackage{algpseudocode}
\PassOptionsToPackage{hidelinks}{hyperref}
\usepackage{hyperref}
\usepackage{orcidlink}
\usepackage[capitalise,noabbrev]{cleveref}
\usepackage{natbib}
\bibliographystyle{plainnat}

\graphicspath{{figures/}}

% Tolerate a partially-populated results directory: a missing table or figure
% becomes a visible placeholder rather than a failed build, so the manuscript
% can be proofread while the protocol is still running.
\newcommand{\inputtable}[1]{%
  \IfFileExists{tables/#1.tex}{\input{tables/#1}}%
  {\begin{table}[t]\centering\fbox{\texttt{\detokenize{#1}}: results pending}\end{table}}}
\newcommand{\resultfig}[2]{%
  \IfFileExists{figures/#1.pdf}{\includegraphics[width=\linewidth]{#1.pdf}}%
  {\fbox{\parbox{0.9\linewidth}{\centering\vspace{2em}\texttt{\detokenize{#1}}: figure pending\vspace{2em}}}}}
% Auto-generated by experiments/make_macros.py -- do not edit.
% Every number quoted in the manuscript prose is defined here and comes
% directly from the JSON payloads written by the experiment scripts.

% Visible placeholder for a value that is still open or obsolete.
\providecommand{\todo}[1]{\textbf{[TODO: #1]}}

\newcommand{\AblChangepointDelay}{311.3}
\newcommand{\AblChangepointMiss}{0.06}
\newcommand{\AblChangepointPfa}{0.000}
\newcommand{\AblDilationDelay}{269.8}
\newcommand{\AblDilationMiss}{0.02}
\newcommand{\AblDilationPfa}{0.000}
\newcommand{\AblDkwDelay}{2025.8}
\newcommand{\AblDkwMiss}{0.68}
\newcommand{\AblDkwPfa}{0.000}
\newcommand{\AblDkwSpeedup}{8.3}
\newcommand{\AblEveryFrameDelay}{3.3}
\newcommand{\AblEveryFrameMiss}{0.00}
\newcommand{\AblEveryFramePfa}{1.000}
\newcommand{\AblFullDelay}{244.5}
\newcommand{\AblFullMiss}{0.01}
\newcommand{\AblFullPfa}{0.000}
\newcommand{\AblLag}{18}
\newcommand{\AblNoCcDelay}{43.1}
\newcommand{\AblNoCcMiss}{0.00}
\newcommand{\AblNoCcPfa}{0.039}
\newcommand{\AblNumInvalid}{5}
\newcommand{\AblNumUnresolved}{14}
\newcommand{\AblNumVariants}{19}
\newcommand{\AblPointPriorDelay}{2800.0}
\newcommand{\AblPointPriorMiss}{1.00}
\newcommand{\AblPointPriorPfa}{0.000}
\newcommand{\AblPooledDelay}{391.0}
\newcommand{\AblPooledMiss}{0.01}
\newcommand{\AblPooledPfa}{0.989}
\newcommand{\AblRawCalDelay}{68.4}
\newcommand{\AblRawCalMiss}{0.00}
\newcommand{\AblRawCalPfa}{0.002}
\newcommand{\AblReps}{200}
\newcommand{\AblationAlpha}{0.010}
\newcommand{\BetaOverDkwSpeedup}{11.9}
\newcommand{\BetaPmin}{0.0037}
\newcommand{\CalDelta}{10^{-3}}
\newcommand{\CalN}{2{,}000}
\newcommand{\CalRawPmin}{1/2001}
\newcommand{\CornerDenseAbsDiff}{0.1}
\newcommand{\CornerDenseCi}{[-0.2, -0.1]}
\newcommand{\CornerDenseCpOnly}{0}
\newcommand{\CornerDenseDiff}{-0.1}
\newcommand{\CornerDenseEpiOnly}{0}
\newcommand{\CornerDenseP}{1.000}
\newcommand{\CornerFloorMiss}{1.00}
\newcommand{\CornerLongGapAbsDiff}{141.1}
\newcommand{\CornerLongGapCi}{[+79.6, +210.2]}
\newcommand{\CornerLongGapCpOnly}{0}
\newcommand{\CornerLongGapDiff}{141.1}
\newcommand{\CornerLongGapEpiOnly}{16}
\newcommand{\CornerLongGapP}{0.000}
\newcommand{\CornerPermanentAbsDiff}{0.1}
\newcommand{\CornerPermanentCi}{[-0.1, -0.0]}
\newcommand{\CornerPermanentCpOnly}{0}
\newcommand{\CornerPermanentDiff}{-0.1}
\newcommand{\CornerPermanentEpiOnly}{0}
\newcommand{\CornerPermanentP}{1.000}
\newcommand{\CornerReps}{200}
\newcommand{\CornerShortGapAbsDiff}{38.5}
\newcommand{\CornerShortGapCi}{[+1.4, +86.5]}
\newcommand{\CornerShortGapCpOnly}{0}
\newcommand{\CornerShortGapDiff}{38.5}
\newcommand{\CornerShortGapEpiOnly}{3}
\newcommand{\CornerShortGapP}{0.250}
\newcommand{\CornerSingleAbsDiff}{55.7}
\newcommand{\CornerSingleCi}{[-111.2, -14.0]}
\newcommand{\CornerSingleCpOnly}{4}
\newcommand{\CornerSingleDiff}{-55.7}
\newcommand{\CornerSingleEpiOnly}{0}
\newcommand{\CornerSingleP}{0.125}
\newcommand{\CostAmortStreamUs}{1.98}
\newcommand{\CostAmortVecUs}{0.15}
\newcommand{\CostBackboneMs}{36.8}
\newcommand{\CostBackboneParams}{1.6}
\newcommand{\CostBetEndToEndUs}{12.2}
\newcommand{\CostBetSteps}{10{,}000}
\newcommand{\CostConformalN}{150{,}000}
\newcommand{\CostConformalUs}{0.34}
\newcommand{\CostFitSec}{0.12}
\newcommand{\CostFrames}{200{,}000}
\newcommand{\CostGnnPerCameraUs}{4.0}
\newcommand{\CostLag}{20}
\newcommand{\CostMonitorSec}{0.122}
\newcommand{\CostRealtimeFactor}{66{,}048}
\newcommand{\CostSentinelPct}{0.0017}
\newcommand{\CostSentinelUs}{0.61}
\newcommand{\CostStepUs}{18.3}
\newcommand{\CostStepVecUs}{2.6}
\newcommand{\CoverDaytimeAbstain}{82}
\newcommand{\CoverDaytimeMiss}{0.90}
\newcommand{\CoverDaytimePfa}{0.000}
\newcommand{\CoverNoRainAbstain}{25}
\newcommand{\CoverNoRainMiss}{0.14}
\newcommand{\CoverNoRainPfa}{0.000}
\newcommand{\CoverRepAbstain}{0}
\newcommand{\CoverRepMiss}{0.01}
\newcommand{\CoverRepPfa}{0.000}
\newcommand{\DelayChangepointmixturePfa}{0.010}
\newcommand{\DelayChangepointmixtureSec}{68.6}
\newcommand{\DelayInadmissibleList}{Fixed threshold, p-value threshold, CUSUM, Shiryaev--Roberts, Parametric e-detector, E-SHIFT}
\newcommand{\DelayNumInadmissible}{6}
\newcommand{\DelaySentinelCi}{65.2}
\newcommand{\DelaySentinelCondSec}{205.6}
\newcommand{\DelaySentinelMiss}{0.015}
\newcommand{\DelaySentinelPfa}{0.000}
\newcommand{\DelaySentinelSec}{244.5}
\newcommand{\DiagAcfRaw}{0.839}
\newcommand{\DiagAcfThin}{0.046}
\newcommand{\DiagLagMax}{21}
\newcommand{\DiagLagMean}{18.8}
\newcommand{\DiagLagMin}{17}
\newcommand{\DiagLjungFrac}{0.183}
\newcommand{\DiagLjungRaw}{$<10^{-6}$}
\newcommand{\DiagLjungThin}{0.005}
\newcommand{\DiagPN}{191{,}890}
\newcommand{\DiagPmin}{0.0018}
\newcommand{\DiagStreams}{60}
\newcommand{\DkwPmin}{0.0441}
\newcommand{\EpiBestCount}{4}
\newcommand{\EpiBestCpDelay}{311}
\newcommand{\EpiBestCpMiss}{0.055}
\newcommand{\EpiBestEpiDelay}{245}
\newcommand{\EpiBestEpiMiss}{0.015}
\newcommand{\EpiBestGainPct}{21}
\newcommand{\EpiBestMissRatio}{3.7}
\newcommand{\EpiEtaMin}{0.0001}
\newcommand{\EpiEventLenSec}{32}
\newcommand{\EpiGridEta}{4}
\newcommand{\EpiGridKappa}{6}
\newcommand{\EpiGridPi}{1}
\newcommand{\EpiGridPoints}{24}
\newcommand{\EpiHardPiCpDelay}{311.3}
\newcommand{\EpiHardPiCpMiss}{0.055}
\newcommand{\EpiHardPiEpiDelay}{244.5}
\newcommand{\EpiHardPiEpiMiss}{0.015}
\newcommand{\EpiManyCpDelay}{242}
\newcommand{\EpiManyCpMiss}{0.010}
\newcommand{\EpiManyCpPfa}{0.000}
\newcommand{\EpiManyEpiDelay}{220}
\newcommand{\EpiManyEpiMiss}{0.000}
\newcommand{\EpiManyEpiPfa}{0.000}
\newcommand{\EpiMaxCount}{8}
\newcommand{\EpiMidCount}{4}
\newcommand{\EpiMidCpDelay}{311}
\newcommand{\EpiMidCpMiss}{0.055}
\newcommand{\EpiMidCpPfa}{0.000}
\newcommand{\EpiMidEpiDelay}{245}
\newcommand{\EpiMidEpiMiss}{0.015}
\newcommand{\EpiMidEpiPfa}{0.000}
\newcommand{\EpiMinPi}{0.35}
\newcommand{\EpiOneCount}{1}
\newcommand{\EpiOneCpDelay}{1202}
\newcommand{\EpiOneCpMiss}{0.425}
\newcommand{\EpiOneCpPfa}{0.000}
\newcommand{\EpiOneEpiDelay}{1257}
\newcommand{\EpiOneEpiMiss}{0.445}
\newcommand{\EpiOneEpiPfa}{0.000}
\newcommand{\EpiPairedCIMany}{[+2.5, +44.3]}
\newcommand{\EpiPairedCIMid}{[+20.4, +120.8]}
\newcommand{\EpiPairedCIOne}{[-111.2, -14.0]}
\newcommand{\EpiPairedCITwo}{[+25.6, +144.8]}
\newcommand{\EpiPairedDiffMany}{22.6}
\newcommand{\EpiPairedDiffMid}{66.8}
\newcommand{\EpiPairedDiffOne}{-55.7}
\newcommand{\EpiPairedDiffTwo}{80.0}
\newcommand{\EpiPairedDiscordantMany}{2}
\newcommand{\EpiPairedDiscordantMid}{8}
\newcommand{\EpiPairedDiscordantOne}{4}
\newcommand{\EpiPairedDiscordantTwo}{6}
\newcommand{\EpiPairedPMcnemarMany}{0.500}
\newcommand{\EpiPairedPMcnemarMid}{0.008}
\newcommand{\EpiPairedPMcnemarOne}{0.125}
\newcommand{\EpiPairedPMcnemarTwo}{0.031}
\newcommand{\EpiPairedPbootMany}{0.026}
\newcommand{\EpiPairedPbootMid}{0.002}
\newcommand{\EpiPairedPbootOne}{0.000}
\newcommand{\EpiPairedPbootTwo}{0.001}
\newcommand{\EpiReps}{200}
\newcommand{\EqCpFifthCiSec}{49.4}
\newcommand{\EqCpFifthDelaySec}{119.6}
\newcommand{\EqCpFifthMiss}{0.010}
\newcommand{\EqCpFifthPfa}{0.000}
\newcommand{\EqCpHalfCiSec}{31.7}
\newcommand{\EqCpHalfDelaySec}{68.6}
\newcommand{\EqCpHalfMiss}{0.005}
\newcommand{\EqCpHalfPfa}{0.010}
\newcommand{\EqCpHundredthCiSec}{92.5}
\newcommand{\EqCpHundredthDelaySec}{311.3}
\newcommand{\EqCpHundredthMiss}{0.055}
\newcommand{\EqCpHundredthPfa}{0.000}
\newcommand{\EqSEFifthCiSec}{49.3}
\newcommand{\EqSEFifthDelaySec}{118.7}
\newcommand{\EqSEFifthMiss}{0.010}
\newcommand{\EqSEFifthPfa}{0.000}
\newcommand{\EqSEHalfCiSec}{31.8}
\newcommand{\EqSEHalfDelaySec}{68.3}
\newcommand{\EqSEHalfMiss}{0.005}
\newcommand{\EqSEHalfPfa}{0.000}
\newcommand{\EqSEHundredthCiSec}{65.2}
\newcommand{\EqSEHundredthDelaySec}{244.5}
\newcommand{\EqSEHundredthMiss}{0.015}
\newcommand{\EqSEHundredthPfa}{0.000}
\newcommand{\ExpNineReps}{200}
\newcommand{\ExpOneHorizonReps}{200}
\newcommand{\ExpOneReps}{1500}
\newcommand{\ExpTenReps}{200}
\newcommand{\ExpThreeReps}{200}
\newcommand{\ExpTwoReps}{200}
\newcommand{\HorizonFixedPfa}{1.000}
\newcommand{\HorizonFixedPfaPct}{100}
\newcommand{\HorizonFixedPfaShort}{0.640}
\newcommand{\HorizonMaxFrames}{240{,}000}
\newcommand{\HorizonMaxHours}{2.7}
\newcommand{\HorizonMinHours}{0.2}
\newcommand{\HorizonPvaluePfa}{1.000}
\newcommand{\HorizonPvaluePfaPct}{100}
\newcommand{\HorizonSentinelPfa}{0.000}
\newcommand{\HorizonWeekBaselineReps}{50}
\newcommand{\HorizonWeekCalDefaultPct}{6.9}
\newcommand{\HorizonWeekCalN}{2{,}160{,}000}
\newcommand{\HorizonWeekFixedOneHour}{0.980}
\newcommand{\HorizonWeekFixedSixHour}{1.000}
\newcommand{\HorizonWeekMaxHours}{168}
\newcommand{\HorizonWeekMaxMFrames}{15.12}
\newcommand{\HorizonWeekPvalueOneHour}{1.000}
\newcommand{\HorizonWeekPvalueSixHour}{1.000}
\newcommand{\HorizonWeekReps}{200}
\newcommand{\HorizonWeekSentinelCiHi}{0.019}
\newcommand{\HorizonWeekSentinelPfaMax}{0.000}
\newcommand{\NetCameras}{12}
\newcommand{\NetCreditWindowSec}{80}
\newcommand{\NetEvalMinutes}{40}
\newcommand{\NetFdrLevel}{0.10}
\newcommand{\NetFleets}{40}
\newcommand{\NetGnnCensDelaySec}{435.3}
\newcommand{\NetGnnDelaySec}{201.1}
\newcommand{\NetGnnFdr}{0.000}
\newcommand{\NetGnnMiss}{0.26}
\newcommand{\NetGnnNullPfa}{0.000}
\newcommand{\NetGnnPower}{0.745}
\newcommand{\NetGnnWindowedCensDelaySec}{462.7}
\newcommand{\NetGnnWindowedDelaySec}{248.6}
\newcommand{\NetGnnWindowedFdr}{0.000}
\newcommand{\NetGnnWindowedMiss}{0.28}
\newcommand{\NetGnnWindowedNullPfa}{0.000}
\newcommand{\NetGnnWindowedPower}{0.724}
\newcommand{\NetHeuristicCensDelaySec}{335.0}
\newcommand{\NetHeuristicDelaySec}{187.7}
\newcommand{\NetHeuristicFdr}{0.000}
\newcommand{\NetHeuristicMiss}{0.17}
\newcommand{\NetHeuristicNullPfa}{0.000}
\newcommand{\NetHeuristicPower}{0.832}
\newcommand{\NetNoGraphCensDelaySec}{343.3}
\newcommand{\NetNoGraphDelaySec}{191.9}
\newcommand{\NetNoGraphFdr}{0.000}
\newcommand{\NetNoGraphMiss}{0.17}
\newcommand{\NetNoGraphNullPfa}{0.000}
\newcommand{\NetNoGraphPower}{0.827}
\newcommand{\NetPropLagSec}{80}
\newcommand{\NetPropProb}{0.45}
\newcommand{\NetTrainFleets}{10}
\newcommand{\PerEpiManyCpMiss}{0.873}
\newcommand{\PerEpiManyCpStreamMiss}{0.000}
\newcommand{\PerEpiManyEpiMiss}{0.873}
\newcommand{\PerEpiManyEpiStreamMiss}{0.000}
\newcommand{\PerEpiMidCpMiss}{0.762}
\newcommand{\PerEpiMidCpStreamMiss}{0.050}
\newcommand{\PerEpiMidEpiMiss}{0.754}
\newcommand{\PerEpiMidEpiStreamMiss}{0.015}
\newcommand{\PerEpiOneCpMiss}{0.440}
\newcommand{\PerEpiOneCpStreamMiss}{0.440}
\newcommand{\PerEpiOneEpiMiss}{0.455}
\newcommand{\PerEpiOneEpiStreamMiss}{0.455}
\newcommand{\QualExampleAvenueAnomalous}{0.096}
\newcommand{\QualExampleAvenueNormal}{0.026}
\newcommand{\QualExamplePedTwoAnomalous}{0.086}
\newcommand{\QualExamplePedTwoNormal}{0.074}
\newcommand{\RFourWarmFlowNormDelayMed}{--}
\newcommand{\RFourWarmFlowNormDetHi}{0.60}
\newcommand{\RFourWarmFlowNormDetLo}{0.00}
\newcommand{\RFourWarmFlowNormDetP}{0.00}
\newcommand{\RFourWarmFlowNormDetk}{0}
\newcommand{\RFourWarmFlowNormDetn}{4}
\newcommand{\RFourWarmFlowNormDynDetk}{0}
\newcommand{\RFourWarmFlowNormDynDetn}{2}
\newcommand{\RFourWarmFlowNormFAhrHi}{48.8}
\newcommand{\RFourWarmFlowNormFAhrLo}{0.2}
\newcommand{\RFourWarmFlowNormFAhrP}{8.8}
\newcommand{\RFourWarmFlowNormFAk}{1}
\newcommand{\RFourWarmFlowNormFAn}{25}
\newcommand{\RFourWarmFlowNormFAp}{0.040}
\newcommand{\RFourWarmFlowNormLag}{5}
\newcommand{\RFourWarmFlowNormStaticDetk}{0}
\newcommand{\RFourWarmFlowNormStaticDetn}{2}
\newcommand{\RFourWarmFlowNullN}{25}
\newcommand{\RFourWarmFlowRawDelayMed}{--}
\newcommand{\RFourWarmFlowRawDetHi}{0.60}
\newcommand{\RFourWarmFlowRawDetLo}{0.00}
\newcommand{\RFourWarmFlowRawDetP}{0.00}
\newcommand{\RFourWarmFlowRawDetk}{0}
\newcommand{\RFourWarmFlowRawDetn}{4}
\newcommand{\RFourWarmFlowRawDynDetk}{0}
\newcommand{\RFourWarmFlowRawDynDetn}{2}
\newcommand{\RFourWarmFlowRawFAhrHi}{102.3}
\newcommand{\RFourWarmFlowRawFAhrLo}{14.2}
\newcommand{\RFourWarmFlowRawFAhrP}{43.8}
\newcommand{\RFourWarmFlowRawFAk}{3}
\newcommand{\RFourWarmFlowRawFAn}{25}
\newcommand{\RFourWarmFlowRawFAp}{0.120}
\newcommand{\RFourWarmFlowRawLag}{7}
\newcommand{\RFourWarmFlowRawStaticDetk}{0}
\newcommand{\RFourWarmFlowRawStaticDetn}{2}
\newcommand{\RFourWarmFlowW}{20}
\newcommand{\RFourWarmPCNormDelayMed}{1.0}
\newcommand{\RFourWarmPCNormDetHi}{0.22}
\newcommand{\RFourWarmPCNormDetLo}{0.00}
\newcommand{\RFourWarmPCNormDetP}{0.04}
\newcommand{\RFourWarmPCNormDetk}{1}
\newcommand{\RFourWarmPCNormDetn}{23}
\newcommand{\RFourWarmPCNormDynDetk}{1}
\newcommand{\RFourWarmPCNormDynDetn}{9}
\newcommand{\RFourWarmPCNormFAhrHi}{62.1}
\newcommand{\RFourWarmPCNormFAhrLo}{17.3}
\newcommand{\RFourWarmPCNormFAhrP}{34.7}
\newcommand{\RFourWarmPCNormFAk}{2}
\newcommand{\RFourWarmPCNormFAn}{45}
\newcommand{\RFourWarmPCNormFAp}{0.044}
\newcommand{\RFourWarmPCNormLag}{5}
\newcommand{\RFourWarmPCNormStaticDetk}{0}
\newcommand{\RFourWarmPCNormStaticDetn}{14}
\newcommand{\RFourWarmPCNullN}{45}
\newcommand{\RFourWarmPCRawDelayMed}{3.5}
\newcommand{\RFourWarmPCRawDetHi}{0.39}
\newcommand{\RFourWarmPCRawDetLo}{0.05}
\newcommand{\RFourWarmPCRawDetP}{0.17}
\newcommand{\RFourWarmPCRawDetk}{4}
\newcommand{\RFourWarmPCRawDetn}{23}
\newcommand{\RFourWarmPCRawDynDetk}{1}
\newcommand{\RFourWarmPCRawDynDetn}{9}
\newcommand{\RFourWarmPCRawFAhrHi}{135.1}
\newcommand{\RFourWarmPCRawFAhrLo}{63.9}
\newcommand{\RFourWarmPCRawFAhrP}{94.7}
\newcommand{\RFourWarmPCRawFAk}{8}
\newcommand{\RFourWarmPCRawFAn}{45}
\newcommand{\RFourWarmPCRawFAp}{0.178}
\newcommand{\RFourWarmPCRawLag}{7}
\newcommand{\RFourWarmPCRawStaticDetk}{3}
\newcommand{\RFourWarmPCRawStaticDetn}{14}
\newcommand{\RFourWarmPCW}{5}
\newcommand{\RFourWarmSupNormDelayMed}{8.3}
\newcommand{\RFourWarmSupNormDetHi}{0.34}
\newcommand{\RFourWarmSupNormDetLo}{0.03}
\newcommand{\RFourWarmSupNormDetP}{0.13}
\newcommand{\RFourWarmSupNormDetk}{3}
\newcommand{\RFourWarmSupNormDetn}{23}
\newcommand{\RFourWarmSupNormDynDetk}{1}
\newcommand{\RFourWarmSupNormDynDetn}{9}
\newcommand{\RFourWarmSupNormFAhrHi}{32.3}
\newcommand{\RFourWarmSupNormFAhrLo}{3.4}
\newcommand{\RFourWarmSupNormFAhrP}{12.6}
\newcommand{\RFourWarmSupNormFAk}{3}
\newcommand{\RFourWarmSupNormFAn}{45}
\newcommand{\RFourWarmSupNormFAp}{0.067}
\newcommand{\RFourWarmSupNormLag}{7}
\newcommand{\RFourWarmSupNormStaticDetk}{2}
\newcommand{\RFourWarmSupNormStaticDetn}{14}
\newcommand{\RFourWarmSupNullN}{45}
\newcommand{\RFourWarmSupRawDelayMed}{2.5}
\newcommand{\RFourWarmSupRawDetHi}{0.93}
\newcommand{\RFourWarmSupRawDetLo}{0.56}
\newcommand{\RFourWarmSupRawDetP}{0.78}
\newcommand{\RFourWarmSupRawDetk}{18}
\newcommand{\RFourWarmSupRawDetn}{23}
\newcommand{\RFourWarmSupRawDynDetk}{7}
\newcommand{\RFourWarmSupRawDynDetn}{9}
\newcommand{\RFourWarmSupRawFAhrHi}{142.6}
\newcommand{\RFourWarmSupRawFAhrLo}{69.1}
\newcommand{\RFourWarmSupRawFAhrP}{101.0}
\newcommand{\RFourWarmSupRawFAk}{6}
\newcommand{\RFourWarmSupRawFAn}{45}
\newcommand{\RFourWarmSupRawFAp}{0.133}
\newcommand{\RFourWarmSupRawLag}{5}
\newcommand{\RFourWarmSupRawStaticDetk}{11}
\newcommand{\RFourWarmSupRawStaticDetn}{14}
\newcommand{\RFourWarmSupW}{5}
\newcommand{\RThreeAvenueFlowAuc}{0.783}
\newcommand{\RThreeAvenueFlowDelayMed}{9.4}
\newcommand{\RThreeAvenueFlowDepCond}{0.203}
\newcommand{\RThreeAvenueFlowDepCondCount}{39}
\newcommand{\RThreeAvenueFlowDepCorr}{0.761}
\newcommand{\RThreeAvenueFlowDepCorrOne}{0.975}
\newcommand{\RThreeAvenueFlowDepSup}{0.000}
\newcommand{\RThreeAvenueFlowDepSupOne}{0.000}
\newcommand{\RThreeAvenueFlowDetHi}{0.11}
\newcommand{\RThreeAvenueFlowDetLo}{0.01}
\newcommand{\RThreeAvenueFlowDetP}{0.04}
\newcommand{\RThreeAvenueFlowDetk}{3}
\newcommand{\RThreeAvenueFlowDetn}{74}
\newcommand{\RThreeAvenueFlowFAHi}{0.206}
\newcommand{\RThreeAvenueFlowFALo}{0.000}
\newcommand{\RThreeAvenueFlowFAhrHi}{31.94}
\newcommand{\RThreeAvenueFlowFAhrLo}{0.00}
\newcommand{\RThreeAvenueFlowFAhrP}{0.00}
\newcommand{\RThreeAvenueFlowFAk}{0}
\newcommand{\RThreeAvenueFlowFAn}{16}
\newcommand{\RThreeAvenueFlowFAp}{0.000}
\newcommand{\RThreeAvenueFlowLag}{20}
\newcommand{\RThreeAvenueFlowLovoK}{1}
\newcommand{\RThreeAvenueFlowLovoN}{16}
\newcommand{\RThreeAvenueFlowNVideos}{21}
\newcommand{\RThreeAvenueFlowWithinLagOne}{0.94}
\newcommand{\RThreeAvenueFlowWithinLagTwentyFive}{0.47}
\newcommand{\RThreeAvenuePCAuc}{0.737}
\newcommand{\RThreeAvenuePCDelayMed}{1.0}
\newcommand{\RThreeAvenuePCDepCond}{0.083}
\newcommand{\RThreeAvenuePCDepCondCount}{20}
\newcommand{\RThreeAvenuePCDepCorr}{0.758}
\newcommand{\RThreeAvenuePCDepCorrOne}{0.987}
\newcommand{\RThreeAvenuePCDepSup}{0.000}
\newcommand{\RThreeAvenuePCDepSupOne}{0.003}
\newcommand{\RThreeAvenuePCDetHi}{0.22}
\newcommand{\RThreeAvenuePCDetLo}{0.06}
\newcommand{\RThreeAvenuePCDetP}{0.12}
\newcommand{\RThreeAvenuePCDetk}{9}
\newcommand{\RThreeAvenuePCDetn}{74}
\newcommand{\RThreeAvenuePCFAHi}{0.302}
\newcommand{\RThreeAvenuePCFALo}{0.002}
\newcommand{\RThreeAvenuePCFAhrHi}{159.22}
\newcommand{\RThreeAvenuePCFAhrLo}{41.52}
\newcommand{\RThreeAvenuePCFAhrP}{86.58}
\newcommand{\RThreeAvenuePCFAk}{1}
\newcommand{\RThreeAvenuePCFAn}{16}
\newcommand{\RThreeAvenuePCFAp}{0.062}
\newcommand{\RThreeAvenuePCLag}{30}
\newcommand{\RThreeAvenuePCLovoK}{0}
\newcommand{\RThreeAvenuePCLovoN}{16}
\newcommand{\RThreeAvenuePCNVideos}{21}
\newcommand{\RThreeAvenuePCWithinLagOne}{0.97}
\newcommand{\RThreeAvenuePCWithinLagTwentyFive}{0.55}
\newcommand{\RThreeBaselineConfigs}{42}
\newcommand{\RThreeBaselineNoKnob}{19}
\newcommand{\RThreeDepCondMax}{0.203}
\newcommand{\RThreeDepCondMin}{0.015}
\newcommand{\RThreeDepCorrMax}{0.985}
\newcommand{\RThreeDepCorrMin}{0.758}
\newcommand{\RThreeDepCorrOneMax}{0.998}
\newcommand{\RThreeDepCountMax}{585}
\newcommand{\RThreeDepCountMin}{20}
\newcommand{\RThreeDepSupMax}{0.032}
\newcommand{\RThreeDepSupOneMax}{0.044}
\newcommand{\RThreeMiviaFlowAuc}{0.344}
\newcommand{\RThreeMiviaFlowDelayMed}{--}
\newcommand{\RThreeMiviaFlowDepCond}{0.104}
\newcommand{\RThreeMiviaFlowDepCondCount}{431}
\newcommand{\RThreeMiviaFlowDepCorr}{0.873}
\newcommand{\RThreeMiviaFlowDepCorrOne}{0.972}
\newcommand{\RThreeMiviaFlowDepSup}{0.000}
\newcommand{\RThreeMiviaFlowDepSupOne}{0.000}
\newcommand{\RThreeMiviaFlowDetHi}{0.07}
\newcommand{\RThreeMiviaFlowDetLo}{0.00}
\newcommand{\RThreeMiviaFlowDetP}{0.00}
\newcommand{\RThreeMiviaFlowDetk}{0}
\newcommand{\RThreeMiviaFlowDetn}{50}
\newcommand{\RThreeMiviaFlowDynDetHi}{0.14}
\newcommand{\RThreeMiviaFlowDynDetLo}{0.00}
\newcommand{\RThreeMiviaFlowDynDetk}{0}
\newcommand{\RThreeMiviaFlowDynDetn}{25}
\newcommand{\RThreeMiviaFlowFAHi}{0.192}
\newcommand{\RThreeMiviaFlowFALo}{0.022}
\newcommand{\RThreeMiviaFlowFAhrHi}{32.22}
\newcommand{\RThreeMiviaFlowFAhrLo}{6.29}
\newcommand{\RThreeMiviaFlowFAhrP}{15.64}
\newcommand{\RThreeMiviaFlowFAk}{4}
\newcommand{\RThreeMiviaFlowFAn}{50}
\newcommand{\RThreeMiviaFlowFAp}{0.080}
\newcommand{\RThreeMiviaFlowLag}{10}
\newcommand{\RThreeMiviaFlowLovoK}{15}
\newcommand{\RThreeMiviaFlowLovoN}{200}
\newcommand{\RThreeMiviaFlowNVideos}{100}
\newcommand{\RThreeMiviaFlowStaticDetHi}{0.14}
\newcommand{\RThreeMiviaFlowStaticDetLo}{0.00}
\newcommand{\RThreeMiviaFlowStaticDetk}{0}
\newcommand{\RThreeMiviaFlowStaticDetn}{25}
\newcommand{\RThreeMiviaFlowWithinLagOne}{0.88}
\newcommand{\RThreeMiviaFlowWithinLagTwentyFive}{0.05}
\newcommand{\RThreeMiviaPCAuc}{0.588}
\newcommand{\RThreeMiviaPCDelayMed}{2.8}
\newcommand{\RThreeMiviaPCDepCond}{0.058}
\newcommand{\RThreeMiviaPCDepCondCount}{585}
\newcommand{\RThreeMiviaPCDepCorr}{0.985}
\newcommand{\RThreeMiviaPCDepCorrOne}{0.998}
\newcommand{\RThreeMiviaPCDepSup}{0.000}
\newcommand{\RThreeMiviaPCDepSupOne}{0.002}
\newcommand{\RThreeMiviaPCDetHi}{0.34}
\newcommand{\RThreeMiviaPCDetLo}{0.10}
\newcommand{\RThreeMiviaPCDetP}{0.20}
\newcommand{\RThreeMiviaPCDetk}{10}
\newcommand{\RThreeMiviaPCDetn}{50}
\newcommand{\RThreeMiviaPCDynDetHi}{0.41}
\newcommand{\RThreeMiviaPCDynDetLo}{0.07}
\newcommand{\RThreeMiviaPCDynDetk}{5}
\newcommand{\RThreeMiviaPCDynDetn}{25}
\newcommand{\RThreeMiviaPCFAHi}{0.314}
\newcommand{\RThreeMiviaPCFALo}{0.086}
\newcommand{\RThreeMiviaPCFAhrHi}{134.52}
\newcommand{\RThreeMiviaPCFAhrLo}{73.33}
\newcommand{\RThreeMiviaPCFAhrP}{100.53}
\newcommand{\RThreeMiviaPCFAk}{9}
\newcommand{\RThreeMiviaPCFAn}{50}
\newcommand{\RThreeMiviaPCFAp}{0.180}
\newcommand{\RThreeMiviaPCLag}{7}
\newcommand{\RThreeMiviaPCLovoK}{20}
\newcommand{\RThreeMiviaPCLovoN}{200}
\newcommand{\RThreeMiviaPCNVideos}{100}
\newcommand{\RThreeMiviaPCStaticDetHi}{0.41}
\newcommand{\RThreeMiviaPCStaticDetLo}{0.07}
\newcommand{\RThreeMiviaPCStaticDetk}{5}
\newcommand{\RThreeMiviaPCStaticDetn}{25}
\newcommand{\RThreeMiviaPCWithinLagOne}{0.92}
\newcommand{\RThreeMiviaPCWithinLagTwentyFive}{0.12}
\newcommand{\RThreeMiviaSupAuc}{0.809}
\newcommand{\RThreeMiviaSupDelayMed}{2.9}
\newcommand{\RThreeMiviaSupDepCond}{0.069}
\newcommand{\RThreeMiviaSupDepCondCount}{583}
\newcommand{\RThreeMiviaSupDepCorr}{0.914}
\newcommand{\RThreeMiviaSupDepCorrOne}{0.984}
\newcommand{\RThreeMiviaSupDepSup}{0.000}
\newcommand{\RThreeMiviaSupDepSupOne}{0.000}
\newcommand{\RThreeMiviaSupDetHi}{0.75}
\newcommand{\RThreeMiviaSupDetLo}{0.47}
\newcommand{\RThreeMiviaSupDetP}{0.62}
\newcommand{\RThreeMiviaSupDetk}{31}
\newcommand{\RThreeMiviaSupDetn}{50}
\newcommand{\RThreeMiviaSupDynDetHi}{0.82}
\newcommand{\RThreeMiviaSupDynDetLo}{0.43}
\newcommand{\RThreeMiviaSupDynDetk}{16}
\newcommand{\RThreeMiviaSupDynDetn}{25}
\newcommand{\RThreeMiviaSupFAHi}{0.192}
\newcommand{\RThreeMiviaSupFALo}{0.022}
\newcommand{\RThreeMiviaSupFAhrHi}{87.76}
\newcommand{\RThreeMiviaSupFAhrLo}{39.75}
\newcommand{\RThreeMiviaSupFAhrP}{60.32}
\newcommand{\RThreeMiviaSupFAk}{4}
\newcommand{\RThreeMiviaSupFAn}{50}
\newcommand{\RThreeMiviaSupFAp}{0.080}
\newcommand{\RThreeMiviaSupLag}{7}
\newcommand{\RThreeMiviaSupLovoK}{19}
\newcommand{\RThreeMiviaSupLovoN}{200}
\newcommand{\RThreeMiviaSupNVideos}{100}
\newcommand{\RThreeMiviaSupStaticDetHi}{0.79}
\newcommand{\RThreeMiviaSupStaticDetLo}{0.39}
\newcommand{\RThreeMiviaSupStaticDetk}{15}
\newcommand{\RThreeMiviaSupStaticDetn}{25}
\newcommand{\RThreeMiviaSupWithinLagOne}{0.91}
\newcommand{\RThreeMiviaSupWithinLagTwentyFive}{0.09}
\newcommand{\RThreePedTwoFlowAuc}{0.763}
\newcommand{\RThreePedTwoFlowDelayMed}{2.7}
\newcommand{\RThreePedTwoFlowDepCond}{0.091}
\newcommand{\RThreePedTwoFlowDepCondCount}{32}
\newcommand{\RThreePedTwoFlowDepCorr}{0.818}
\newcommand{\RThreePedTwoFlowDepCorrOne}{0.946}
\newcommand{\RThreePedTwoFlowDepSup}{0.000}
\newcommand{\RThreePedTwoFlowDepSupOne}{0.004}
\newcommand{\RThreePedTwoFlowDetHi}{1.00}
\newcommand{\RThreePedTwoFlowDetLo}{0.62}
\newcommand{\RThreePedTwoFlowDetP}{0.92}
\newcommand{\RThreePedTwoFlowDetk}{11}
\newcommand{\RThreePedTwoFlowDetn}{12}
\newcommand{\RThreePedTwoFlowFAHi}{0.906}
\newcommand{\RThreePedTwoFlowFALo}{0.008}
\newcommand{\RThreePedTwoFlowFAhrHi}{1586.86}
\newcommand{\RThreePedTwoFlowFAhrLo}{7.21}
\newcommand{\RThreePedTwoFlowFAhrP}{284.81}
\newcommand{\RThreePedTwoFlowFAk}{1}
\newcommand{\RThreePedTwoFlowFAn}{3}
\newcommand{\RThreePedTwoFlowFAp}{0.333}
\newcommand{\RThreePedTwoFlowLag}{5}
\newcommand{\RThreePedTwoFlowLovoK}{1}
\newcommand{\RThreePedTwoFlowLovoN}{16}
\newcommand{\RThreePedTwoFlowNVideos}{12}
\newcommand{\RThreePedTwoFlowWithinLagOne}{0.82}
\newcommand{\RThreePedTwoFlowWithinLagTwentyFive}{0.16}
\newcommand{\RThreePedTwoPCAuc}{0.828}
\newcommand{\RThreePedTwoPCDelayMed}{2.2}
\newcommand{\RThreePedTwoPCDepCond}{0.015}
\newcommand{\RThreePedTwoPCDepCondCount}{23}
\newcommand{\RThreePedTwoPCDepCorr}{0.839}
\newcommand{\RThreePedTwoPCDepCorrOne}{0.965}
\newcommand{\RThreePedTwoPCDepSup}{0.032}
\newcommand{\RThreePedTwoPCDepSupOne}{0.044}
\newcommand{\RThreePedTwoPCDetHi}{0.85}
\newcommand{\RThreePedTwoPCDetLo}{0.28}
\newcommand{\RThreePedTwoPCDetP}{0.58}
\newcommand{\RThreePedTwoPCDetk}{7}
\newcommand{\RThreePedTwoPCDetn}{12}
\newcommand{\RThreePedTwoPCFAHi}{0.708}
\newcommand{\RThreePedTwoPCFALo}{0.000}
\newcommand{\RThreePedTwoPCFAhrHi}{1050.63}
\newcommand{\RThreePedTwoPCFAhrLo}{0.00}
\newcommand{\RThreePedTwoPCFAhrP}{0.00}
\newcommand{\RThreePedTwoPCFAk}{0}
\newcommand{\RThreePedTwoPCFAn}{3}
\newcommand{\RThreePedTwoPCFAp}{0.000}
\newcommand{\RThreePedTwoPCLag}{5}
\newcommand{\RThreePedTwoPCLovoK}{1}
\newcommand{\RThreePedTwoPCLovoN}{16}
\newcommand{\RThreePedTwoPCNVideos}{12}
\newcommand{\RThreePedTwoPCWithinLagOne}{0.87}
\newcommand{\RThreePedTwoPCWithinLagTwentyFive}{-0.00}
\newcommand{\RThreeSupCUSUMDetP}{0.50}
\newcommand{\RThreeSupCUSUMDetk}{25}
\newcommand{\RThreeSupCUSUMDetn}{50}
\newcommand{\RThreeSupCUSUMFAhrP}{42.4}
\newcommand{\RThreeSupCUSUMFAk}{5}
\newcommand{\RThreeSupCUSUMFAn}{50}
\newcommand{\RThreeSupEShiftDetP}{0.76}
\newcommand{\RThreeSupEShiftDetk}{38}
\newcommand{\RThreeSupEShiftDetn}{50}
\newcommand{\RThreeSupEShiftFAhrP}{183.2}
\newcommand{\RThreeSupEShiftFAk}{8}
\newcommand{\RThreeSupEShiftFAn}{50}
\newcommand{\RThreeSupFixedDetP}{0.52}
\newcommand{\RThreeSupFixedDetk}{26}
\newcommand{\RThreeSupFixedDetn}{50}
\newcommand{\RThreeSupFixedFAhrP}{109.5}
\newcommand{\RThreeSupFixedFAk}{2}
\newcommand{\RThreeSupFixedFAn}{50}
\newcommand{\RThreeSupPValueDetP}{0.62}
\newcommand{\RThreeSupPValueDetk}{31}
\newcommand{\RThreeSupPValueDetn}{50}
\newcommand{\RThreeSupPValueFAhrP}{201.1}
\newcommand{\RThreeSupPValueFAk}{4}
\newcommand{\RThreeSupPValueFAn}{50}
\newcommand{\RThreeSupParamDetP}{0.70}
\newcommand{\RThreeSupParamDetk}{35}
\newcommand{\RThreeSupParamDetn}{50}
\newcommand{\RThreeSupParamFAhrP}{78.2}
\newcommand{\RThreeSupParamFAk}{7}
\newcommand{\RThreeSupParamFAn}{50}
\newcommand{\RThreeSupSEDetP}{0.62}
\newcommand{\RThreeSupSEDetk}{31}
\newcommand{\RThreeSupSEDetn}{50}
\newcommand{\RThreeSupSEFAhrP}{60.3}
\newcommand{\RThreeSupSEFAk}{4}
\newcommand{\RThreeSupSEFAn}{50}
\newcommand{\RThreeSupSRDetP}{0.52}
\newcommand{\RThreeSupSRDetk}{26}
\newcommand{\RThreeSupSRDetn}{50}
\newcommand{\RThreeSupSRFAhrP}{49.1}
\newcommand{\RThreeSupSRFAk}{6}
\newcommand{\RThreeSupSRFAn}{50}
\newcommand{\RThreeWithinLagOneMax}{0.97}
\newcommand{\RThreeWithinLagOneMin}{0.82}
\newcommand{\RThreeWithinLagTwentyFiveMax}{0.55}
\newcommand{\RThreeWithinLagTwentyFiveMin}{-0.00}
\newcommand{\RealAvenueAuc}{0.783}
\newcommand{\RealAvenueCiHiHeadline}{\todo{obsolete}}
\newcommand{\RealAvenueEpiMax}{328}
\newcommand{\RealAvenueEpiMedian}{29}
\newcommand{\RealAvenueEpiMin}{2}
\newcommand{\RealAvenueFullLag}{207}
\newcommand{\RealAvenueLag}{\todo{obsolete}}
\newcommand{\RealAvenueMissHeadline}{\todo{obsolete}}
\newcommand{\RealAvenueNEpisodes}{74}
\newcommand{\RealAvenueNTestVideos}{21}
\newcommand{\RealAvenueNTrainVideos}{16}
\newcommand{\RealAvenuePfaHeadline}{\todo{obsolete}}
\newcommand{\RealAvenueShiftFrac}{8.9}
\newcommand{\RealAvenueWorstAlpha}{\todo{obsolete}}
\newcommand{\RealAvenueWorstPfa}{\todo{obsolete}}
\newcommand{\RealDelayAvenueMiss}{\todo{obsolete}}
\newcommand{\RealDelayAvenueNBaselines}{\todo{obsolete}}
\newcommand{\RealDelayAvenueNBaselinesPass}{\todo{obsolete}}
\newcommand{\RealDelayAvenuePfa}{\todo{obsolete}}
\newcommand{\RealDelayAvenueSec}{\todo{obsolete}}
\newcommand{\RealDelayPedTwoMiss}{\todo{obsolete}}
\newcommand{\RealDelayPedTwoNBaselines}{\todo{obsolete}}
\newcommand{\RealDelayPedTwoNBaselinesPass}{\todo{obsolete}}
\newcommand{\RealDelayPedTwoPfa}{\todo{obsolete}}
\newcommand{\RealDelayPedTwoSec}{\todo{obsolete}}
\newcommand{\RealDelayReps}{200}
\newcommand{\RealLagReps}{\todo{obsolete}}
\newcommand{\RealPedTwoAuc}{0.763}
\newcommand{\RealPedTwoCiHiHeadline}{\todo{obsolete}}
\newcommand{\RealPedTwoEpiMax}{180}
\newcommand{\RealPedTwoEpiMedian}{140}
\newcommand{\RealPedTwoEpiMin}{86}
\newcommand{\RealPedTwoFullLag}{357}
\newcommand{\RealPedTwoLag}{\todo{obsolete}}
\newcommand{\RealPedTwoMissHeadline}{\todo{obsolete}}
\newcommand{\RealPedTwoNEpisodes}{12}
\newcommand{\RealPedTwoNTestVideos}{12}
\newcommand{\RealPedTwoNTrainVideos}{16}
\newcommand{\RealPedTwoPfaHeadline}{\todo{obsolete}}
\newcommand{\RealPedTwoShiftFrac}{32.6}
\newcommand{\RealPedTwoWorstAlpha}{\todo{obsolete}}
\newcommand{\RealPedTwoWorstPfa}{\todo{obsolete}}
\newcommand{\RealReps}{400}
\newcommand{\ShiftBreakResidualMiss}{0.26}
\newcommand{\ShiftBreakResidualPfa}{0.070}
\newcommand{\ShiftBreakResidualPfaHi}{0.114}
\newcommand{\ShiftBreakResidualPfaLo}{0.042}
\newcommand{\ShiftBreakWeightedMiss}{1.00}
\newcommand{\ShiftBreakWeightedPfa}{0.000}
\newcommand{\ShiftDefGateMinBlocks}{0.99}
\newcommand{\ShiftDefGateMinFrac}{1.00}
\newcommand{\ShiftDefMissDelta}{0.070}
\newcommand{\ShiftDefMissDeltaMax}{0.05}
\newcommand{\ShiftDefNullStreams}{1500}
\newcommand{\ShiftDefOneHalfResCiHi}{0.006}
\newcommand{\ShiftDefOneHalfResCiLo}{0.001}
\newcommand{\ShiftDefOneHalfResMiss}{0.060}
\newcommand{\ShiftDefOneHalfResPfa}{0.002}
\newcommand{\ShiftDefOneHalfResVerdict}{valid}
\newcommand{\ShiftDefOneHalfWtMiss}{0.540}
\newcommand{\ShiftDefOneHalfWtPfa}{0.000}
\newcommand{\ShiftDefOneResCiHi}{0.003}
\newcommand{\ShiftDefOneResCiLo}{0.000}
\newcommand{\ShiftDefOneResMiss}{0.025}
\newcommand{\ShiftDefOneResPfa}{0.000}
\newcommand{\ShiftDefOneResVerdict}{valid}
\newcommand{\ShiftDefOneWtMiss}{0.095}
\newcommand{\ShiftDefOneWtPfa}{0.000}
\newcommand{\ShiftDefRecommendWeighted}{keep unweighted}
\newcommand{\ShiftDefSignalStreams}{200}
\newcommand{\ShiftDefThreeResCiHi}{0.075}
\newcommand{\ShiftDefThreeResCiLo}{0.051}
\newcommand{\ShiftDefThreeResMiss}{0.265}
\newcommand{\ShiftDefThreeResPfa}{0.062}
\newcommand{\ShiftDefThreeResVerdict}{invalid}
\newcommand{\ShiftDefThreeWtMiss}{1.000}
\newcommand{\ShiftDefThreeWtPfa}{0.000}
\newcommand{\ShiftDefTwoResCiHi}{0.021}
\newcommand{\ShiftDefTwoResCiLo}{0.009}
\newcommand{\ShiftDefTwoResMiss}{0.115}
\newcommand{\ShiftDefTwoResPfa}{0.014}
\newcommand{\ShiftDefTwoResVerdict}{unresolved}
\newcommand{\ShiftDefTwoWtMiss}{0.930}
\newcommand{\ShiftDefTwoWtPfa}{0.000}
\newcommand{\ShiftMaxRatio}{3.0}
\newcommand{\ShiftMondrianDelay}{14.2}
\newcommand{\ShiftMondrianMiss}{0.00}
\newcommand{\ShiftMondrianPfa}{1.000}
\newcommand{\ShiftPooledDelay}{8.0}
\newcommand{\ShiftPooledMiss}{0.00}
\newcommand{\ShiftPooledPfa}{1.000}
\newcommand{\ShiftResidualBreaks}{3.0}
\newcommand{\ShiftResidualDelay}{956.3}
\newcommand{\ShiftResidualLRDelay}{2800.0}
\newcommand{\ShiftResidualLRMiss}{1.00}
\newcommand{\ShiftResidualLRPfa}{0.000}
\newcommand{\ShiftResidualMiss}{0.26}
\newcommand{\ShiftResidualPfa}{0.070}
\newcommand{\ShiftResidualSafeUpTo}{--}
\newcommand{\ShiftResidualUnresolvedRatios}{$\times$1.0, $\times$1.5, $\times$2.0}
\newcommand{\SimCalN}{150{,}000}
\newcommand{\SimDiurnalFrames}{40{,}000}
\newcommand{\SimDiurnalMin}{26.7}
\newcommand{\SpeedupChangepointmixture}{0.3}
\newcommand{\ValidityAlphaLoose}{0.20}
\newcommand{\ValidityMaxSentinelPfa}{0.001}
\newcommand{\ValidityPfaOracleLoose}{0.044}
\newcommand{\ValidityPfaPvalue}{0.998}
\newcommand{\ValidityPfaPvaluePct}{99.8}
\newcommand{\ValidityPfaSentinel}{0.000}
\newcommand{\ValidityPvalueArl}{11{,}155}
\newcommand{\ValidityPvalueArlMin}{7.4}
\newcommand{\ValidityTightestAlpha}{0.005}
\newcommand{\ValidityTightestCiHi}{0.003}
\newcommand{\ValidityTightestPfa}{0.000}
\newcommand{\ValidityWorstCiHi}{0.004}
\newcommand{\WorkedClimbSec}{23}
\newcommand{\WorkedDelaySec}{23.4}
\newcommand{\WorkedEp}{[320, 352]\,s, [392, 424]\,s, [455, 487]\,s}
\newcommand{\WorkedNEpisodes}{3}
\newcommand{\WorkedT}{1200}


\theoremstyle{plain}
\newtheorem{theorem}{Theorem}
\newtheorem{lemma}[theorem]{Lemma}
\newtheorem{proposition}[theorem]{Proposition}
\newtheorem{corollary}[theorem]{Corollary}
\theoremstyle{definition}
\newtheorem{definition}{Definition}
\newtheorem{remark}{Remark}

\newcommand{\Pinf}{\mathbb{P}_{\infty}}
\newcommand{\Einf}{\mathbb{E}_{\infty}}
\newcommand{\Fcal}{\mathcal{F}}
\newcommand{\Gfilt}{\mathcal{G}}
\newcommand{\cD}{\mathcal{D}}
\newcommand{\ind}[1]{\mathbf{1}\{#1\}}
\newcommand{\ARL}{\mathrm{ARL}_0}
\newcommand{\ADD}{\mathrm{ADD}}
\newcommand{\PFA}{\mathrm{PFA}}

\title{\textbf{SENTINEL-E}: Anytime-Valid Sequential Detection for\\
       Streaming Visual Anomaly Surveillance}

\author{}
\date{}

\hypersetup{pdfauthor={},pdftitle={SENTINEL-E},pdfsubject={},pdfkeywords={}}

\begin{document}
\maketitle
\thispagestyle{empty}
\begin{center}
  \emph{Manuscript prepared for double-anonymous peer review. Author names,
  affiliations, and declarations appear on the separate title page.}
\end{center}
\vspace{0.5em}

\begin{abstract}
\noindent
Visual-anomaly detectors, such as illegal-dumping detectors, are evaluated per clip but deployed on continuous streams, where alerting at the first threshold crossing is optional stopping. Fed valid 1\% p-values, a per-frame rule raises a false alarm on \ValidityPfaPvaluePct\% of 40-minute null streams. We introduce the episodic e-process, a mixture over episodes that start, flicker, end and recur. It is a non-negative supermartingale when the p-values are conditionally super-uniform, so Ville's inequality gives time-uniform false-alarm control; it contains the changepoint mixture as the case where episodes never end. A two-state recursion evaluates it in constant time per frame. In simulation it raised no false alarm on 200 week-long streams; per-frame thresholds reach 0.980 within an hour. Paired comparisons show episodes shorten delay for long, intermittent, recurring events, while the changepoint mixture is faster for one episode. On real video, operating points come from training videos and each test video is monitored alone. Thinned p-values remain serially dependent, so real-video validity is verified empirically. Per-video false-alarm proportions are consistent with $\alpha=0.1$ on all three corpora (3, 16 and 50 null videos). With a Mivia-IWDD-trained dumping detector, SENTINEL-E detects 31 of 50 events and raises alarms on 4 of 50 normal videos at $\alpha=0.1$; among methods with at most 5 such videos, only the p-value threshold matches it, with 3.3 times more alarms per hour. Detection power depends on the backbone. On the short pedestrian videos, several baselines detect more events with more false alarms. SENTINEL-E costs under two microseconds per frame.
\end{abstract}

\vspace{0.5em}
\noindent\textbf{Keywords:} video anomaly detection; illegal dumping; anytime-valid inference; e-values; conformal prediction; sequential change detection; video surveillance.

\section{Introduction}
\label{sec:intro}

Illegal dumping is a persistent municipal problem with an obvious
computer-vision framing. The field has adopted that framing consistently: label
clips as containing a dumping event or not, train a classifier and report
clip-level accuracy. The 2026 Illegal Waste Dumping Detection (IWDD) contest \citep{bouwmans2026overview} gave this
protocol a shared benchmark. Ten teams trained on Mivia-IWDD (400 balanced
videos with onset timestamps) and were ranked on a private test set of 100 videos.
Seven entries are published
\citep{almorsy2026x3d,kim2026kfold,scarrica2026garberus,kairanbay2026vlm,zbuce2026novelty,bhat2026imslab,putzu2026unicass}.
Where they count false alarms, they count them per clip, as in ``six false alarms
on thirty negative videos'' \citep{almorsy2026x3d}. That is specificity, not a
rate over a monitoring period. The larger Mivia-IWDD-500 release
\citep{greco2026iwdd} postdates the contest.

The framing does not match deployment. A camera does not return a verdict on a
curated clip. It produces frames for weeks, and the system alerts when its
evidence first crosses a threshold. The municipality manages the number of wasted
enforcement dispatches over a season. Testing a stream repeatedly and stopping at
the first exceedance is optional stopping, and it invalidates the level at which
the threshold was calibrated.

The gap is larger than it looks. A per-frame rule that false-alarms with
probability $q$ raises at least one alarm over $T$ frames with probability
$1-(1-q)^{T}$, which tends to one for any $q>0$. Feed that rule p-values that are
valid at the $1\%$ level. The probability of at least one false alarm on a $40$-minute null stream, the false-alarm probability (PFA), is then \ValidityPfaPvaluePct\% (\cref{fig:validity}\,(a)). A
raw threshold tuned to a $10^{-3}$ per-frame rate reaches a false-alarm probability
of $0.980$ within one hour and $1.000$ from six hours (\cref{tab:horizon_week}). A
clip-level evaluation never runs the detector for hours, so it cannot see this.

Nothing here is specific to dumping. Loitering, wrong-way motion or an abandoned
object faces the same gap. The method therefore targets streaming visual anomaly
surveillance, with dumping as the motivating case. Real-data validation covers UCSD
Ped2 and CUHK Avenue, which are pedestrian-walkway corpora, and Mivia-IWDD-500,
which contains illegal-dumping footage. Each corpus is run with an optical-flow and
a PatchCore backbone, and each test video is monitored alone. The Mivia labels give
onset times only, and the real streams are short. The long-horizon claims therefore
rest on simulation (\cref{sec:limitations}).

\subsection{Why the existing sequential machinery needs an episode axis}

The tools for optional stopping are mature. Ville's inequality \citep{ville1939}
bounds the probability that a non-negative supermartingale ever exceeds a level.
Betting against the null turns the accumulated wealth into a time-uniform
certificate \citep{shafer2019game,ramdas2023savi}. Conformal p-values
\citep{vovk2005alrw} supply the bets. E-detectors \citep{shin2024edetectors} adapt
the machinery to change detection, and e-BH \citep{wang2022ebh} combines evidence
across streams. We know of no application of this machinery to waste, dumping or
pollution surveillance.

Importing it unchanged would not be enough. Every one of these procedures models a
persistent change. At an unknown frame $\nu$ the law switches from $P_0$ to $P_1$
and stays there, and the mixtures run over $\nu$ alone.

An illegal dumping event is an \emph{episode}. A van pulls up, the offender unloads
for forty seconds and drives away. Within the episode the signal is intermittent:
the offender is occluded, leaves the frame and returns. Then the stream reverts to
null. A site used once is often used again weeks later.

A persistent-change model pays for this mismatch. A changepoint-mixture wealth
process keeps betting at full aggression after the episode ends. It also counts the
quiet frames inside the episode against the alternative. Both effects lengthen
delay, and the second grows with intermittency. In paired simulations on grids of equal
size, the episodic mixture reduces the mean censored delay by 141.1\,s (95\% CI +79.6
to +210.2\,s) for four recurring 1600-frame episodes with $\pi=0.2$ and a mean gap of
12{,}000 frames (\cref{tab:corner_sweep}). A single episode reverses the sign, by 55.7\,s (95\% CI 14.0 to
111.2\,s) in favour of the changepoint mixture (\cref{tab:corner_sweep}). \cref{sec:episodes} maps the regimes.

\subsection{Contributions}

The paper makes three kinds of contribution: theory, an algorithm and an
evaluation method.

\begin{itemize}
\item \textbf{Theory: an episodic e-process.} Mixing over episodes, not
      changepoints, gives an exact non-negative supermartingale, so Ville's
      inequality applies unchanged (\cref{thm:episodic}). It equals the changepoint
      mixture exactly when the episode-end rate is zero (\cref{prop:generalisation}).

\item \textbf{Theory: exact conformal levels.} Each attainable p-value level is
      replaced by its exact Beta order-statistic bound (\cref{thm:beta}). At
      $n=\CalN$ and $\delta=\CalDelta$, the smallest usable p-value falls from
      $\DkwPmin$ to $\BetaPmin$, an improvement of $\BetaOverDkwSpeedup\times$ over
      the Dvoretzky--Kiefer--Wolfowitz bound.

\item \textbf{Algorithm: a constant-time recursion.} A forward recursion runs the
      mixture in $O(1)$ time and memory per frame (\cref{eq:forward,alg:episodic}).
      Direct evaluation over active starts costs $O(t)$. The guarantee layer adds
      under two microseconds per frame (\cref{sec:compute}).

\item \textbf{A learned prior and a posterior.} A learned
      spatial prior enters only as a predictable modulation, so the guarantee holds
      whatever it outputs (\cref{thm:predictable}). The posterior probability of
      being inside an episode is bounded in $[0,1]$ and is available to future fleet
      controllers. In the tested fleets, detection power comes from the e-BH
      combination, and the learned prior lengthens delay (\cref{sec:network}).

\item \textbf{An evaluation of guarantees.} We report realised false-alarm rates
      over up to \ExpOneReps{} streams per operating point, against six baselines,
      including E-SHIFT \citep{khan2026eshift}. Delay is reported censored as well as
      conditionally (\cref{sec:problem}). Real video is monitored per video, with
      operating points chosen on training videos only (\cref{sec:real-data}).
      Dilated betting is evaluated as a design option (\cref{sec:dilation}).
\end{itemize}

\paragraph{What we do not claim.}
The changepoint mixture with retained tail mass is not new. It is the
PFA-controlled e-detector of \citet{saha2026nonpartitioned}, which gives
$\mathbb{P}(\tau<\infty)\le\alpha$. Our contribution is the episode axis on top of
it, and the $O(1)$ recursion that makes that axis cheap. Conformal prediction and
e-BH are existing machinery. \cref{sec:limitations} sets out what the evaluation
does and does not establish.

\section{Related work}
\label{sec:related}

\paragraph{Waste and dumping detection.}
TACO \citep{proenca2020taco}, ZeroWaste \citep{bashkirova2022zerowaste},
UAVVaste \citep{kraft2021uavvaste}, AerialWaste \citep{torres2023aerialwaste},
MARIDA \citep{kikaki2022marida} and MADOS \citep{kikaki2024mados} are image,
aerial, satellite or industrial datasets. None is a continuous video stream.

The 2026 IWDD contest \citep{bouwmans2026overview} ranked ten teams on Mivia-IWDD by
precision, recall, $F_1$, normalised notification delay, frame rate and peak memory.
Seven entries are published: 3D CNNs with temporal localisation
\citep{almorsy2026x3d}, transformer ensembles \citep{kim2026kfold}, multi-head
classifiers \citep{scarrica2026garberus}, vision--language reasoning
\citep{kairanbay2026vlm}, novelty detection \citep{zbuce2026novelty},
onset-centric estimation \citep{bhat2026imslab} and a two-stage recurrent framework
\citep{putzu2026unicass}. SENTINEL-E consumes a frozen backbone's scores, so it
complements these entries. None of them reports a false-alarm rate over a
continuous horizon, an average run length or a delay under continuous monitoring.
The contest's ROC analysis is also per clip. Our real-data protocol uses the larger
Mivia-IWDD-500 release \citep{greco2026iwdd}.

\paragraph{Video anomaly detection benchmarks.}
CUHK Avenue \citep{lu2013avenue} and UCSD Ped2 \citep{mahadevan2010ucsd} are
classical benchmarks, and recent work reports frame-level area under the ROC curve (AUC) on them. One
reconstruction-based method reports 99.03\% AUC on Ped2 and 92.40\% on Avenue
\citep{dong2025rvadign}. For a survey, see \citet{liu2024generalized}. AUC is
computed after the fact on a labelled test set. It does not measure how often a
rule alarms on an unlabelled stream, which is what our protocol measures.
\citet{rashidi2026benchmarkauc} show the consequence. Off-the-shelf feature sets from
four backbones, evaluated on four datasets, average 70.4\% AUC on their training
distribution and fall to 49.9\%, which is chance, under cross-dataset shift. At an
operating point that reaches 0.9 recall on the source dataset, the false-alarm
rate is around $3.2\times 10^{4}$ per hour.

Our real-data backbones are optical flow, estimated by Farneb\"{a}ck's method
\citep{farneback2003}, and PatchCore. PatchCore scores each patch by its
nearest-neighbour distance to a memory bank of frozen ImageNet features from a
residual network \citep{roth2022total,he2016deep}.

\paragraph{Online detection with a false-alarm bound.}
The closest prior work is \citet{doshi2021online}. It gives an online detector for
surveillance video with an asymptotic bound on the false-alarm rate. Our guarantee
is non-asymptotic and time-uniform, so it holds at all monitoring times at once.
Doshi and Yilmaz also treat video anomaly detection as continual learning
\citep{doshi2022rethinking}.

\paragraph{Sequential detection of changes that end.}
Quickest-detection theory mostly treats permanent shifts, but a sub-literature
treats changes of finite duration. Transient change detection
\citep{guepie2012transient,guepie2017transient} minimises worst-case missed
detection under a false-alarm constraint. Intermittent changes of unknown duration
are treated in \citet{sokolov2023intermittent}. \citet{fuh2015hmm} use a two-state
hidden Markov model, the closest classical relative of the chain in
\cref{eq:episodic}. Epidemic changepoint methods \citep{fisch2022capa,juodakis2020epidemic}
handle a segment that departs from and returns to a background level. These methods
use CUSUM, SPRT or window statistics with asymptotic or fixed-sample guarantees.
None is time-uniform. We add an anytime-valid, conformal-fed construction of the
finite-duration change. The finite-duration change itself is not new, and
\cref{sec:episodic} states that scope.

Mixing over the trajectories of a two-state chain belongs to the fixed-share, or
switching-expert, family \citep{herbster1998fixedshare,vanerven2012switch}. That
literature does not give the martingale property under optional stopping, which
\cref{thm:episodic} establishes. Closer to our architecture are the conformal test
martingales of \citet{vovk2003exchangeability}. They pair a conformal front end with
a Ville-thresholded martingale for change and exchangeability testing. We know of no
e-process or anytime-valid conformal method applied to waste, dumping or pollution
surveillance.

\paragraph{Sequential change detection.}
CUSUM \citep{page1954}, Shiryaev's statistic \citep{shiryaev1963} and
Shiryaev--Roberts \citep{roberts1966} are the classical procedures, with optimality
due to \citet{lorden1971} and \citet{moustakides1986}. A full treatment is in
\citet{tartakovsky2014sequential}. Their guarantees are asymptotic
statements about the average run length (ARL) under a correctly specified pre-change law. Footage
scores are autocorrelated, heavy-tailed and diurnally non-stationary, and both
procedures lose control in our experiments. \citet{shin2024edetectors} give a
nonparametric e-detector framework with non-asymptotic ARL bounds. \citet{saha2026nonpartitioned}
extend it to composite pre-change classes. Their PFA-controlled variant mixes over
start times with weights summing to one, which gives $\mathbb{P}(\tau<\infty)\le\alpha$.
\cref{prop:generalisation} reduces to that variant at zero episode-end rate.

All of these mix over the changepoint location of a persistent shift. None adds an
episode-end axis with a time-uniform, optional-stopping-safe guarantee. Cost also
separates them. Evaluated exactly over all active starts, a non-partitioned
e-detector costs $O(t)$ per observation \citep{saha2026nonpartitioned}. Under a
geometric prior over start times, the sum telescopes to the $O(1)$ recursion of
\cref{eq:forward} at $\eta=0$. The episodic process keeps that cost while adding the
episode-end axis.

\paragraph{Anytime-valid inference and e-values.}
Testing by betting \citep{shafer2021testing} turns Ville's inequality into a design
principle, and \citet{waudby2024betting} develop the betting strategies we adapt.
\citet{wang2022ebh} introduce e-BH, which controls the false discovery rate (FDR) under
arbitrary dependence and underpins our fleet layer. The closest work in spirit is
E-SHIFT \citep{khan2026eshift}, a sub-exponential e-process for distribution shift
in streaming systems. We reimplement it as a baseline. It differs in three ways.
Its log-moment-generating-function bound is estimated, so by the authors' account
validity is asymptotic in the calibration size. It assumes i.i.d.\ calibration and
deployment scores, which is stronger than the exchangeability our conformal layer
needs, and a diurnal cycle breaks it. Its wealth restarts after each alert, so
evidence arriving hours into a stream must first repay what was burned earlier. Our
restart convention between alarms has the same effect.

\paragraph{Conformal prediction.}
Split conformal prediction \citep{vovk2005alrw,angelopoulos2023gentle} supplies the
distribution-free p-values we bet on. Conformal p-values also test for outliers
\citep{bates2023testing}, and exchangeability can be relaxed for sequential data
\citep{barber2023beyond,xu2023sequential}. \citet{vovk2012conditional} studies
calibration-conditional validity, which our Beta levels address. Covariate shift is
handled by likelihood-ratio weighting \citep{tibshirani2019covariate}, which we
implement and evaluate. Conformalised regression residuals \citep{lei2018distribution}
pursue conditional coverage by other means. Layer 1 adds the sequential requirement.
The p-values must be super-uniform conditionally on the past, not merely marginally,
or the wealth process is not a supermartingale.

\paragraph{Control charts and ROC practice.}
Control charts are tuned to an in-control average run length, and their limits are
revisited when the process changes \citep{montgomery2020introduction}. Our
time-uniform false-alarm probability is the horizon-free analogue: it bounds the
chance of any alarm at any time. ROC analysis fixes a per-decision operating point
\citep{fawcett2006introduction}. A budget of false alarms per monitoring period
depends on how many decisions are made, which a per-decision point does not show.
\cref{sec:limitations} returns to alarm budgets.

\section{Method}
\label{sec:method}

\subsection{Problem formulation}
\label{sec:problem}

A camera produces frames $x_1, x_2, \ldots$, and a frozen detector $g$ produces scores $s_t = g(x_t) \in \mathbb{R}$, with higher meaning more dumping-like. Each frame also yields a low-dimensional context vector $c_t \in \mathbb{R}^d$: time of day, a weather flag, ambient luminance and a background-subtraction motion measure. Any deployed camera can observe these without privileged knowledge. The operator supplies a per-camera calibration set $\cD = \{(s_i^{\mathrm{cal}}, c_i^{\mathrm{cal}})\}_{i=1}^{n}$ of confirmed no-dumping frames.

Write $\Pinf$ for the law of the stream under ``no dumping ever'' and $\mathbb{P}_{\nu}$ for the law when dumping begins at frame $\nu$. A detection rule is a stopping time $\tau$ with respect to the observation filtration $(\Fcal_t)$. We require:

\begin{definition}[Time-uniform false-alarm control]
\label{def:tuv}
A rule has time-uniform false-alarm level $\alpha$ if
$\Pinf(\tau < \infty) \le \alpha$.
\end{definition}

The quantifier is the whole point. A per-frame guarantee $\Pinf(\text{alarm at } t) \le \alpha$ for each fixed $t$ says nothing about $\Pinf(\exists t : \text{alarm at } t)$, and the gap grows with the horizon. \Cref{def:tuv} is what a municipality buys when told a detector has a $1\%$ false-alarm rate.

Dumping is episodic, so we write $\mathbb{P}_{\mathcal{E}}$ for the law when a set of episodes $\mathcal{E} = \{[\nu_1, \nu_1 + L_1], \ldots\}$ occurs, within each of which a fraction $\pi$ of frames are anomalous. The single persistent change is the special case of one episode with $L = \infty$ and $\pi = 1$.

We report, over Monte Carlo replicates: $\PFA$, the false-alarm probability: the probability that a null stream of a stated length raises at least one alarm; $\ARL$, the average run length: the mean frames to a false alarm on a null stream, censored at the horizon; $\ADD$, the average detection delay: the mean frames from the first onset $\nu_1$ to the first alarm at or after it, over runs that alarm; and the miss rate, the fraction of post-change segments ending unflagged.

\paragraph{A measurement pitfall, and the fix.}
Conditional $\ADD$ is averaged over a different subset of runs for each method. A detector that misses the hard events is scored only on the easy ones, and so looks fast. Because the methods compared here differ in miss rate, we also report a censored delay, in which a run that never alarms is charged the whole remaining horizon. The two can rank methods in opposite order. In \cref{sec:episodes} the changepoint mixture has the shorter conditional delay and the longer censored delay. Where they disagree, the censored figure is the comparable one, and head-to-head claims in this paper rest on it.

Every comparison uses a method's \emph{realised} $\PFA$, not its nominal level. A method that misreports its own error rate would otherwise appear to dominate.

\paragraph{Real-video protocol.}
Each real video is monitored as its own stream, with a fresh detector. Calibration uses training videos only. The betting lag and every baseline knob are selected by leave-one-video-out on training videos, so no test frame informs an operating point (\cref{sec:real-data}). The splicing utility \texttt{build\_stream\_from\_clips} remains in the library for constructing long synthetic streams from clips.

\subsection{Layer 1: conformal calibration}
\label{sec:layer1}

\paragraph{Betting rate.}
The e-process of \cref{sec:episodic} needs increments that are conditionally super-uniform given the past. Video violates this badly. Four lags are in play, and they are distinct:
\begin{itemize}
\item \emph{Pooled autocorrelation length.} On concatenated training videos the autocorrelation stays above $0.05$ beyond 200 lags. This reflects differences between video means, so we do not use it.
\item \emph{Within-video autocorrelation.} At lag one it is $0.82$--$0.97$ on the real backbones (\cref{sec:real-data}) and $\DiagAcfRaw$ in simulation.
\item \emph{Decorrelation lag $\ell$.} In simulation it is estimated from the calibration set alone: $\ell = \AblLag$.
\item \emph{Betting lag $L^\ast$.} On real footage it is chosen by leave-one-video-out on training videos, because short, heterogeneous videos make an autocorrelation estimate unreliable.
\end{itemize}
For the simulated streams, $\ell$ is estimated after projecting out the context-predictable drift with a quadratic least-squares fit, since the slow diurnal component would otherwise return a uselessly large $\ell$. We place one bet per $\ell$ frames. With $\ell = \AblLag$, bets occur more than once per second, far finer than a dumping event. Calibration frames are consecutive video too, so they are thinned by the same lag. The effective calibration size is $n/\ell$ (or $n/L^\ast$ on real footage), and the guarantee is stated in that number. \Cref{sec:ablation} shows that betting on every frame is the largest single source of invalidity we measured.

\paragraph{Conditional calibration.}
Let $S$ be a nonconformity score. The split-conformal p-value of a deployment frame is $p = (1 + |\{i : S_i \ge S\}|)/(n+1)$, which is super-uniform under exchangeability. A night-time frame is not exchangeable with a daylight-dominated $\cD$, so we conformalise the context-normalised residual
\begin{equation}
\label{eq:residual}
  z \;=\; \frac{s - \hat g(c)}{\hat\sigma(c)},
  \qquad
  \hat g(c) \approx \mathbb{E}[s \mid c],
  \quad
  \hat\sigma(c) \approx \mathbb{E}\bigl[|s - \hat g(c)| \,\big|\, c\bigr].
\end{equation}
Here $\hat g$ and $\hat\sigma$ are ridge regressions on a quadratic-with-interactions basis in $c$, fitted on one half of $\cD$ and conformalised on the other. The split makes the conformalisation residuals independent of the fitted models, so \cref{thm:beta} applies conditionally on them. The remaining assumption is that the normalised residual is independent of $c$. The normalisation provides this, not the split (\cref{sec:limitations}). Mondrian calibration is also implemented, since it is exactly group-conditional. \cref{sec:shift} shows where each wins.

Because a regression can be extrapolated but a guarantee cannot, frames whose context falls outside the region covered by $\cD$ are reported inactive, and the detector bets neutrally on them. The abstention rate reports how often this happens. Under shift the recommended safeguard is the stricter support region, with the Mahalanobis slack set to 0.5 instead of the default 1.5, and the abstention rate is the monitor that signals recalibration (\cref{sec:shift}).

\paragraph{Calibration-conditional validity.}
The same finite $\cD$ scores every frame for months, so the guarantee must hold conditionally on $\cD$. Conditionally, the p-values deviate from uniform by $O(n^{-1/2})$. The standard remedy adds the Dvoretzky--Kiefer--Wolfowitz deviation $\varepsilon_n(\delta) = \sqrt{\log(2/\delta)/(2n)}$ \citep{dkw1956,massart1990} to every p-value. At $n=\CalN{}$, $\delta=\CalDelta{}$ this moves the smallest attainable p-value from $\CalRawPmin{}$ to $\DkwPmin{}$, which is too blunt. A betting detector uses only the tail, so we use the exact distribution instead.

\begin{theorem}[Exact calibration-conditional levels]
\label{thm:beta}
Let $S_1,\dots,S_n$ be i.i.d.\ from a continuous $F$ and let $S$ be an
independent draw from $F$. For $j = 0,\dots,n-1$ let $\delta_j > 0$ with
$\sum_j \delta_j \le \delta$, and set
\begin{equation}
\label{eq:betalevels}
  q_j \;=\; \mathrm{Beta}^{-1}\bigl(1-\delta_j;\; j+1,\; n-j\bigr),
  \qquad q_n = 1 .
\end{equation}
Define $\tilde p = q_{J}$ with $J = |\{i : S_i \ge S\}|$. Then there is an event
$E$ with $\mathbb{P}(E) \ge 1-\delta$, measurable with respect to
$S_1,\dots,S_n$, on which
$\mathbb{P}\bigl(\tilde p \le u \mid S_1,\dots,S_n\bigr) \le u$
for every $u \in [0,1]$.
\end{theorem}

\begin{proof}
Conditionally on the calibration sample, $\{J \le j\} = \{S > S_{(n-j)}\}$ where
$S_{(1)} \le \cdots \le S_{(n)}$ are the order statistics, so
$\mathbb{P}(J \le j \mid \cD) = 1 - F(S_{(n-j)})$. Since $F$ is continuous,
$F(S_{(k)}) \sim \mathrm{Beta}(k, n+1-k)$, hence
$1 - F(S_{(n-j)}) \sim \mathrm{Beta}(j+1, n-j)$. By \cref{eq:betalevels},
$\mathbb{P}\bigl(1 - F(S_{(n-j)}) > q_j\bigr) = \delta_j$. Let $E$ be the
intersection over $j$ of the complements; a union bound gives
$\mathbb{P}(E) \ge 1 - \sum_j \delta_j \ge 1-\delta$. On $E$,
$\mathbb{P}(\tilde p \le q_j \mid \cD) = \mathbb{P}(J \le j \mid \cD) \le q_j$
for every $j$, and since $\tilde p$ takes values only in $\{q_0,\dots,q_n\}$
and $q_j$ is non-decreasing in $j$, super-uniformity at the atoms extends to all
$u$.
\end{proof}

We spend the budget as $\delta_j \propto (j+1)^{-2}$, concentrating confidence where the betting functions look. Under this allocation $\{q_j\}$ is verified non-decreasing in $j$, which identifies $\{\tilde p \le q_j\} = \{J \le j\}$ in the proof above. The implementation takes a running maximum as a safeguard for any allocation where this fails. At $n=\CalN{}$, $\delta=\CalDelta{}$ this gives $q_0 = \BetaPmin{}$, against DKW's $\DkwPmin{}$. \Cref{thm:beta} assumes i.i.d.\ calibration draws, which is why the calibration set is thinned. At 1,500 null streams, unthinned calibration stayed valid (\cref{sec:ablation}), so the experiments do not show that thinning is needed in practice.

When deployment reweights context that $\cD$ covers, weighted conformal calibration \citep{tibshirani2019covariate} applies. We use a logistic density-ratio estimate, gated behind a two-sample test, and replace $n$ in \cref{eq:betalevels} by the Kish effective sample size \citep{kish1965}. This substitution has no calibration-conditional proof, because \cref{thm:beta} assumes equal weights. The weighted path is therefore a heuristic, evaluated empirically in \cref{sec:shift}, and it is not the recommended safeguard.

\subsection{Layer 2: the episodic e-process}
\label{sec:episodic}

Let $p_1, p_2, \ldots$ be the p-values at the betting instants. A \emph{betting function} is a non-increasing $f : (0,1] \to [0,\infty)$ with $\int_0^1 f \le 1$. Then $\mathbb{E}_{H_0}[f(p)] \le 1$ for any super-uniform $p$, so $f(p)$ is an e-value. We use the power family $f_\kappa(p) = \kappa p^{\kappa-1}$, $\kappa \in (0,1]$, with $\kappa = 1$ the neutral bet $f \equiv 1$.

\subsubsection{Mixing over episodes}

Let $z_t \in \{\text{quiet}, \text{active}\}$ be a two-state Markov chain with onset probability $\mathbb{P}(\text{quiet} \to \text{active}) = \rho$ and \emph{end} probability $\mathbb{P}(\text{active} \to \text{quiet}) = \eta$, started quiet. The wealth mixes over every trajectory of that chain:
\begin{equation}
\label{eq:episodic}
  W_t \;=\; \sum_{z_{1:t}} \mathbb{P}(z_{1:t})
            \prod_{s \le t} f(p_s)^{\ind{z_s = \text{active}}} .
\end{equation}
Each trajectory contributes the product of betting factors over the frames it declares active, weighted by its prior probability. The classical changepoint mixture is the special case in which the chain, once active, never returns. It mixes over when a change starts. \Cref{eq:episodic} mixes over when episodes start and stop, and over how many occur.

The sum has $2^t$ terms and collapses to two numbers. Writing $A_t$ and $Q_t$ for the parts of the sum whose trajectories end active and quiet respectively,
\begin{equation}
\label{eq:forward}
  A_t = \bigl[(1-\eta)A_{t-1} + \rho\,Q_{t-1}\bigr] f(p_t),
  \qquad
  Q_t = \eta\,A_{t-1} + (1-\rho)\,Q_{t-1},
\end{equation}
with $A_0 = 0$, $Q_0 = 1$ and $W_t = A_t + Q_t$. This is the forward recursion of a hidden Markov model. It costs $O(1)$ per frame in time and memory.

\begin{theorem}[Time-uniform validity of the episodic e-process]
\label{thm:episodic}
Under $\Pinf$, if the $p_t$ are conditionally super-uniform given
$\Fcal_{t-1}$, then $(W_t)$ defined by \cref{eq:episodic} is a non-negative
supermartingale with $W_0 = 1$, and for any $\alpha \in (0,1)$
\begin{equation*}
  \Pinf\bigl(\exists\, t \ge 1 : W_t \ge 1/\alpha\bigr) \;\le\; \alpha .
\end{equation*}
\end{theorem}

\begin{proof}
Non-negativity is immediate and $W_0 = Q_0 = 1$. From \cref{eq:forward},
$A_{t}$ carries exactly one factor $f(p_t)$ and $Q_t$ carries none, so
\begin{align*}
  \mathbb{E}[W_t \mid \Fcal_{t-1}]
   &= \bigl[(1-\eta)A_{t-1} + \rho Q_{t-1}\bigr]\,
      \mathbb{E}[f(p_t) \mid \Fcal_{t-1}]
      \;+\; \eta A_{t-1} + (1-\rho) Q_{t-1} \\
   &\le (1-\eta)A_{t-1} + \rho Q_{t-1} + \eta A_{t-1} + (1-\rho) Q_{t-1}
    \;=\; A_{t-1} + Q_{t-1} \;=\; W_{t-1},
\end{align*}
using $\mathbb{E}[f(p_t) \mid \Fcal_{t-1}] \le 1$ and the fact that the
transition probabilities out of each state sum to one. Ville's inequality
\citep{ville1939} applied to the non-negative supermartingale $(W_t)$ started at
one gives the bound.
\end{proof}

\begin{corollary}[End-to-end level]
\label{cor:composition}
Let $E$ be the calibration event of \cref{thm:beta}, with
$\mathbb{P}(E) \ge 1 - \delta$, on which the deployment p-values $\tilde p_t$
are conditionally super-uniform given $\Fcal_{t-1}$. Then the episodic
e-process run on $(\tilde p_t)$ satisfies
$\Pinf(\exists\, t \ge 1 : W_t \ge 1/\alpha) \le \alpha + \delta$.
\end{corollary}
\begin{proof}
On $E$, \cref{thm:episodic} applies verbatim and gives probability at most
$\alpha$ conditional on $\cD$; off $E$, the probability is bounded by $1$, and
$\mathbb{P}(E^c) \le \delta$. Removing the conditioning,
$\Pinf(\exists t : W_t \ge 1/\alpha) \le \mathbb{P}(E)\cdot\alpha +
\mathbb{P}(E^c)\cdot 1 \le \alpha + \delta$.
\end{proof}
The end-to-end level is therefore $\alpha + \delta$, as the library documents in \texttt{sentinel\_e/pipeline.py}. With the default $\delta = 10^{-3}$, the relative loosening is $20\%$ at the smallest level in \cref{sec:validity}, $\alpha = 0.005$, and it falls as $\alpha$ grows. Every level plotted in this paper should be read as $\alpha + \delta$.

The episode structure costs nothing in the guarantee, because the transition probabilities telescope exactly. A term-by-term mixture over candidate starts costs $O(t)$ work per observation \citep{saha2026nonpartitioned}, whereas \cref{eq:forward} keeps two numbers per grid point. The geometric-prior changepoint mixture also reduces to $O(1)$ (\cref{prop:generalisation}). The gain is therefore not speed over that special case, but episodes of finite duration at the same cost (\cref{sec:compute}).

\begin{proposition}[Strict generalisation]
\label{prop:generalisation}
With $\eta = 0$ the recursion \cref{eq:forward} coincides with the mixture over
changepoint locations under a geometric prior with hazard $\rho$, with the
unspent tail mass retained. Discarding $Q_t$ yields a statistic that is not a
supermartingale and that tends to the classical Shiryaev--Roberts statistic
only in the limit $\rho \to 0$ after rescaling.
\end{proposition}

\begin{proof}
With $\eta = 0$, $Q_t = (1-\rho)Q_{t-1}$ so $Q_t = (1-\rho)^t$ is the prior mass
on ``no episode has begun by $t$'', and
$A_t = (A_{t-1} + \rho Q_{t-1}) f(p_t)$ with $\rho Q_{t-1} = w_t$ the geometric
prior mass on a change at $t$. Unrolling gives
$A_t = \sum_{j \le t} w_j \prod_{s=j}^{t} f(p_s)$, whence
$W_t = \sum_{j\le t} w_j \prod_{s=j}^t f(p_s) + \sum_{j>t} w_j$, Shiryaev's
Bayesian statistic under the geometric prior with hazard $\rho$
\citep{shiryaev1963}. Dropping $Q_t$ leaves $\sum_{j\le t} w_j \prod_{s=j}^t
f(p_s)$ alone. This is not a supermartingale: the telescoping in
\cref{thm:episodic} fails by exactly $w_t$ per step. Taking the
$\rho \to 0$ limit under the appropriate rescaling recovers the classical
Shiryaev--Roberts sum with unit start weights \citep{roberts1966}.
\end{proof}

The statement thus separates two objects that are easy to conflate: the changepoint mixture, which the episodic process contains exactly at $\eta=0$, and Shiryaev--Roberts, which differs from it at any fixed $\rho>0$.

\subsubsection{Intermittency by dilation (a negative result)}
\label{sec:dilation}

Within an episode only a fraction $\pi$ of frames show the act, which suggests \emph{dilating} the emission:
\begin{equation}
\label{eq:dilated}
  f_\pi(p) \;=\; (1-\pi) + \pi\, f(p).
\end{equation}
This is a legal betting function. It is a convex combination of the neutral bet and $f$, hence non-increasing with $\int_0^1 f_\pi \le 1$, so \cref{thm:episodic} applies verbatim. It does not help. Mixing over $\kappa$ already covers intermittency, because a milder bet suits an episode with fewer anomalous frames. In the ablation, dilation lengthens censored detection delay from \AblFullDelay\,s to \AblDilationDelay\,s (Online Resource~1, Table~S7), and no intermittency level repays that cost. The default grid therefore sets $\pi = 1$. The construction stays in the implementation for backbones whose anomalous frames are much stronger relative to quiet ones than ours.

\subsubsection{Staying parameter-free}
\label{sec:parameterfree}

Nothing above needs $(\rho, \eta, \pi, \kappa)$. We place a prior grid over all four and run one recursion per grid point. The mixture $W_t = \sum_g \lambda_g W_t^{(g)}$ with $\sum_g \lambda_g = 1$ is a convex combination of supermartingales, hence a supermartingale. The deployed grid is \EpiGridKappa{} values of $\kappa$ by \EpiGridEta{} of $\eta$ by \EpiGridPi{} of $\pi$ (dilation off by default), giving \EpiGridPoints{} points. The cost is a constant number of operations per grid point, still $O(1)$ in $t$. An operator does not know how long a dumping event lasts, and the grid spans everything from a ten-second burst to a permanent change.

\subsubsection{The episode posterior}

The recursion also yields, at no extra cost, the posterior probability that the camera is inside an episode now,
\begin{equation}
\label{eq:posterior}
  \Pi_t \;=\; \frac{A_t}{A_t + Q_t} \;\in\; [0, 1],
\end{equation}
prior-averaged over the grid. It is bounded and interpretable as a probability, and it is comparable between a camera installed last week and one that has run for a month. \Cref{sec:layer3} uses it as the message the camera graph exchanges.

\subsubsection{Numerics}

Working in log space would need two \texttt{logaddexp} calls per grid point per frame. Instead we renormalise $(A_t, Q_t)$ to sum to one at each step and accumulate the log of the scale factor, the standard HMM scaling trick. The two-state simplex cannot underflow. \Cref{alg:episodic} gives the per-frame update, which transcribes the inner loop of \texttt{sentinel\_e/episodic.py}.

\begin{algorithm}[t]
\caption{One frame of the episodic e-process, vectorised over the prior grid}
\label{alg:episodic}
\begin{algorithmic}[1]
\Require prior grid $\{(\kappa_g, \eta_g, \pi_g)\}_{g=1}^G$ with log-prior
  weights $\log\lambda_g$, $\sum_g \lambda_g = 1$; onset hazard $\rho$ (or
  $\rho_t$, if predictable modulation is enabled, \cref{thm:predictable});
  alarm threshold $1/\alpha$
\State $A_g \gets 0,\; Q_g \gets 1,\; s_g \gets 0 \;\; \forall g$
  \Comment{$s_g$: per-grid-point log-scale; $A_g + Q_g = 1$ is the invariant}
\For{each betting instant $t = 1, 2, \ldots$}
  \State receive conformal p-value $p_t$ (or skip this $t$: thinning / abstention)
  \For{each grid point $g$ \textbf{(vectorised, not a real loop)}}
    \State $f_g \gets (1 - \pi_g) + \pi_g\, \kappa_g\, p_t^{\,\kappa_g - 1}$
      \Comment{dilated power bet, \cref{eq:dilated}; $\pi_g = 1$ if dilation off}
    \State $A_g' \gets \bigl[(1-\eta_g)\, A_g + \rho\, Q_g\bigr] f_g$
    \State $Q_g' \gets \eta_g\, A_g + (1-\rho)\, Q_g$
      \Comment{\cref{eq:forward}, one grid point}
    \State $c_g \gets A_g' + Q_g'$
      \Comment{restores the invariant; $c_g > 0$ always}
    \State $A_g \gets A_g' / c_g,\;\; Q_g \gets Q_g' / c_g$
      \Comment{HMM scaling trick: the two-state simplex cannot underflow}
    \State $s_g \gets s_g + \log c_g$
  \EndFor
  \State $\log W_t \gets \operatorname{logsumexp}_g\bigl(\log\lambda_g + s_g\bigr)$
    \Comment{mixture happens here, in log space, not by renormalising a joint sum}
  \State $\Pi_t \gets \sum_g \lambda_g\, e^{s_g} A_g \,\big/\, e^{\log W_t}$
    \Comment{episode posterior, \cref{eq:posterior}}
  \If{$\log W_t \ge \log(1/\alpha)$}
    \State \textbf{alarm}; if restarting, reset $A_g \gets 0,\, Q_g \gets 1,\,
      s_g \gets 0\;\; \forall g$
  \EndIf
\EndFor
\end{algorithmic}
\end{algorithm}

Two points about \cref{alg:episodic} are easy to miss from \cref{eq:forward} alone. Each grid point keeps its own invariant $A_g + Q_g = 1$ and log-scale $s_g$. The mixture is a single \texttt{logsumexp}, evaluated only when $\log W_t$ or $\Pi_t$ is read out, so adding grid points lengthens a vector without changing the algorithm. The normalising division is exact, and the only error beyond rounding is what the oracle p-value arm of \cref{sec:validity} is designed to expose.

\paragraph{A worked numeric trace.}
\Cref{tab:worked-recursion} runs four steps of \cref{alg:episodic} at a single grid point ($\kappa = 0.3$, $\eta = 0.1$, $\rho = 0.01$, dilation off), so every entry can be checked with a calculator. Here $c_t = A_t' + Q_t'$ and $W_t = \exp(\sum_{s\le t} \log c_s)$.

\begin{table}[t]
\centering
\caption{Four steps of \cref{eq:forward} at a single grid point
($A_0 = 0$, $Q_0 = 1$), computed to the precision shown; a reader can
reproduce every column from the one before it and the stated $p_t$.}
\label{tab:worked-recursion}
\begin{tabular}{rrrrrrrrr}
\toprule
$t$ & $p_t$ & $f_t$ & $A_t'$ & $Q_t'$ & $c_t$ & $A_t$ & $Q_t$ & $W_t$ \\
\midrule
1 & 0.02 & 4.6387 & 0.046387 & 0.990000 & 1.036387 & 0.044759 & 0.955241 & 1.036387 \\
2 & 0.40 & 0.5697 & 0.028393 & 0.950165 & 0.978558 & 0.029015 & 0.970985 & 1.014165 \\
3 & 0.90 & 0.3230 & 0.011570 & 0.964176 & 0.975746 & 0.011857 & 0.988143 & 0.989568 \\
4 & 0.05 & 2.4425 & 0.050202 & 0.979447 & 1.029649 & 0.048756 & 0.951244 & 1.018907 \\
\bottomrule
\end{tabular}
\end{table}

A small p-value at $t=1$ ($p_1 = 0.02$) multiplies $A_t$ by $f_1 \approx 4.64$, yet $W_1$ rises by under $4\%$, because the prior mass $Q_0 = 1$ on ``no episode yet'' dilutes the bet. At $t=3$, $W_t$ falls below one. Wealth need not grow, only have conditional expectation at most its previous value. Over four frames $A_t \ll Q_t$, so $\Pi_4 \approx 0.049$ reflects too little evidence, not a fault.

\subsection{Layer 3: the camera network}
\label{sec:layer3}

\paragraph{A predictable spatial prior.}
Offenders return to the same lay-by and, when disturbed, move a few hundred metres and dump again. An episode at one camera is therefore prior information about its neighbours. We build a graph over the fleet, with edge weights combining spatial proximity and static-context similarity. A two-layer message-passing network maps each camera's features to an onset hazard $\rho_{i,t}$ and a stake scale $m_{i,t}$.

The features use the episode posterior of \cref{eq:posterior}. Raw log-wealth grows without bound and depends on installation time, so a neighbourhood average of it would be dominated by the oldest camera. $\Pi_{i,t}$ is bounded and answers the same question at every camera, so a neighbourhood mean is meaningful. The network sees $\Pi_{i,t-1}$, the normalised log-wealth, three exponential moving averages of $-\log p$, static context and degree, all measurable at $t-1$.

Because the network passes messages, $\rho_{i,t}$ and $m_{i,t}$ depend on the \emph{fleet's} joint history. Write $\Gfilt_{t-1} = \sigma(p_{j,s} : j \in \text{fleet},\, s \le t-1)$ for that filtration. Then $\Fcal_{t-1} \subseteq \Gfilt_{t-1}$ is camera $i$'s own, and the theorem below is stated at the joint level.

\begin{theorem}[Learned control cannot break the guarantee]
\label{thm:predictable}
Suppose $\rho_{i,t} \in [0,1]$ and $m_{i,t} \in [0,1]$ are
$\Gfilt_{t-1}$-measurable, the map $(m,\kappa) \mapsto f_{\pi,\kappa_{\mathrm{eff}}}$
with $\kappa_{\mathrm{eff}} = 1-m(1-\kappa)$ is jointly measurable, and $p_{i,t}$
is conditionally super-uniform given $\Gfilt_{t-1}$ (the joint history, not
merely camera $i$'s own). Then \cref{thm:episodic} holds verbatim for every
camera, for any such choice of $\rho_{i,t}, m_{i,t}$.
\end{theorem}

\begin{proof}
For $m_{i,t} \in [0,1]$ and $\kappa \in (0,1]$ we have
$\kappa_{\mathrm{eff}} \in (0,1]$, so $f_{\kappa_{\mathrm{eff}}}$ is a betting
function and by \cref{eq:dilated} so is its dilation, for every fixed value of
$\kappa_{\mathrm{eff}}$. Since $(m_{i,t}, \pi)$ is $\Gfilt_{t-1}$-measurable and
the map to $f_{\pi,\kappa_{\mathrm{eff}}}$ is jointly measurable in its
argument and in $p_{i,t}$, conditioning on $\Gfilt_{t-1}$ freezes
$\kappa_{\mathrm{eff}}$ to its realised value and
$\mathbb{E}[f_{\pi,\kappa_{\mathrm{eff}}}(p_{i,t}) \mid \Gfilt_{t-1}] \le 1$
follows from conditional super-uniformity at that fixed value. This is the
standard freezing argument for a regular conditional distribution, not a
step that treats $f_{\kappa_{\mathrm{eff}}}$ as linear in $\kappa_{\mathrm{eff}}$.
Replacing $\rho$ by $\rho_{i,t}$ in \cref{eq:forward} leaves the transition
probabilities out of each state summing to one, which is all the telescoping in
\cref{thm:episodic} uses, now run against $\Gfilt_{t-1}$ throughout. The
supermartingale step is therefore unchanged.
\end{proof}

The hypothesis that matters in practice is conditional super-uniformity given the \emph{joint} fleet history. This is stronger than \cref{sec:layer1} establishes for one camera, because the message-passing layer couples cameras that share weather, lighting and traffic. \Cref{sec:diagnostics} does not certify it, so the theorem is conditional on it. Unconditionally, the learned outputs enter only as a predictable reweighting of a supermartingale step. A badly behaved component therefore cannot cause an invalid step that a well-behaved one would not also cause. \Cref{tab:network}'s all-null columns check the practical consequence.

Training is offline on held-out simulated fleets. It maximises post-change log-wealth growth, the Kelly/GROW criterion \citep{kelly1956}, with a penalty on wealth burned at unaffected cameras. Growth rate proxies delay, since delay is roughly the time for wealth to climb to $1/\alpha$ at that rate. The recursion is differentiable, so gradients reach both heads.

\paragraph{Fleet decisions.}
At a stopping time $\tau_i = \min(\text{alarm}, T)$, optional stopping gives $\Einf[W_{i,\tau_i}] \le 1$, so each camera's wealth is a valid e-value. We apply e-BH \citep{wang2022ebh}: sort $E_{(1)} \ge \cdots \ge E_{(K)}$ and reject the top $k^\star = \max\{k : E_{(k)} \ge K/(\alpha k)\}$. This controls FDR at $\alpha$ under \emph{arbitrary} dependence, which the setting needs because adjacent cameras share weather, lighting and traffic. BH on p-values would need positive regression dependence or a Benjamini--Yekutieli $\log K$ penalty \citep{benjamini2001}. e-BH needs neither. The wealth is read at a stopping time, since a running maximum of the wealth would not be a valid e-value.

\section{Experiments and results}
\label{sec:results}

\subsection{Protocol}
\label{sec:protocol}

Two sources of evidence answer different questions. Simulated streams vary the
horizon, the episode structure and the calibration shift, and each run
reproduces from its seed. Real video (\cref{sec:real-data}) tests the guarantee
on public corpora with real detector scores, but each corpus is short and
covers one backbone at a time. On real video, the betting lag and every
baseline knob are chosen on training videos only, by leave-one-video-out. Each
test video is monitored alone, and no test frame is reused.

\paragraph{Simulated streams.}
SENTINEL-E consumes one scalar per frame, so the simulator generates that
scalar. The null score has a diurnal cycle and a two-state weather chain, and
both shift its distribution after calibration. Innovations are Student-$t$, and
the scores are autocorrelated. A log-scale Ornstein--Uhlenbeck process drives
the backbone's false-positive rate, and it is absent from the operator's
clock-and-weather context. Dumping events are episodic: a set number of episodes
per camera, separated by exponential gaps, with a fraction $\pi$ of frames
anomalous inside each episode. Episodes propagate along the camera graph with a
lag.

The simulated day has \SimDiurnalFrames{} frames, which is \SimDiurnalMin{}\,min
at 25\,fps. This is the diurnal period in Experiments 1--12. Their streams last
40 to 60 minutes. A shorter horizon sweep, up to \HorizonMaxHours\,h, is in Online Resource~1,
Sect.~S3. Only the week-long experiment in \cref{sec:validity} uses a full
24\,h simulated day.

\paragraph{Baselines.}
A fixed score threshold, and a per-frame threshold on the same conformal p-values
that SENTINEL-E consumes, which isolates the stopping rule from calibration. The
others are CUSUM \citep{page1954}, Shiryaev--Roberts \citep{roberts1966}, a
parametric e-detector with a Gaussian likelihood-ratio bet, a reimplementation of
E-SHIFT \citep{khan2026eshift}, and the changepoint mixture. The changepoint
mixture is an instance of the non-partitioned e-detector of
\citet{saha2026nonpartitioned} by \cref{prop:generalisation}, so \cref{sec:episodes}
compares against a published anytime-valid method under the same conformal front
end. All methods share one calibration set.


Every method restarts after an alarm, as an operator does once an alert has been
investigated. Under the single-shot convention of \cref{def:tuv}, a detector
stops at its first alarm and records a miss for every later event. Restarting
leaves the null metrics unchanged, because the false-alarm probability and the
time to the first alarm do not depend on what follows.


Monte Carlo rates carry Wilson 95\% intervals. An upper limit of zero counts
gives the strongest statement that a zero count supports. Each experiment draws
its own seed range, so the same configuration in two tables is an independent
replicate. The paired comparisons in \cref{sec:episodes} share seeds within each
comparison. Every quantity quoted in the text is generated into
\texttt{results/tables/macros.tex} by \texttt{experiments/make\_macros.py}, which
reads the saved experiment output. Timings are single-thread CPU.

\subsection{Real-video validation}
\label{sec:real-data}

Real video tests whether the guarantee holds per video on footage that the
simulator does not produce. Three corpora supply it: CUHK Avenue
\citep{lu2013avenue}, UCSD Ped2 \citep{mahadevan2010ucsd} and Mivia-IWDD
\citep{greco2026iwdd}, an illegal-dumping corpus. Each corpus is run with two
frozen backbones, and no backbone is fine-tuned.

\paragraph{Data and backbones.} Ped2 has \RealPedTwoNTrainVideos{} training
videos, all normal, and \RealPedTwoNTestVideos{} test videos from one fixed
camera. Avenue has \RealAvenueNTrainVideos{} training and
\RealAvenueNTestVideos{} test videos from one fixed camera. Mivia-IWDD has 200
training videos, all labelled negative, and 100 test videos: 50 negative, 25
static dumping and 25 dynamic dumping. Each Mivia video comes from a different
camera, so Mivia calibration pools cameras. The labels mark only the
dumping onset. Mivia is sampled at 10\,fps; the other two corpora at 25\,fps.

The flow backbone is Farneback optical-flow magnitude \citep{farneback2003},
averaged per frame at $64\times64$ resolution and median-filtered over three
frames within each video. The appearance backbone is PatchCore
\citep{roth2022total} with a frozen WideResNet-50-2 pretrained on ImageNet.
It uses layers 2 and 3, a 1\% coreset of training patches, and the maximum
over patches of the nearest-neighbour distance as the frame score. Its frame
AUCs are \RThreePedTwoFlowAuc{} (Ped2 flow), \RThreePedTwoPCAuc{} (Ped2 PatchCore),
\RThreeAvenueFlowAuc{} (Avenue flow), \RThreeAvenuePCAuc{} (Avenue PatchCore),
\RThreeMiviaFlowAuc{} (Mivia flow) and \RThreeMiviaPCAuc{} (Mivia PatchCore).
The Mivia flow score is below chance. PatchCore is a strong training-free appearance
backbone, but it is not the state of the art on these benchmarks. Its frame AUC is
0.828 on Ped2, 0.737 on Avenue (below flow's 0.783) and 0.588 on Mivia. The power
results below therefore bound what such frozen backbones deliver, not what the best
published video anomaly detectors would deliver. We did not add a state-of-the-art video anomaly detector. Such models are trained per dataset, and evaluating them without leakage would need its own training-only selection protocol.

\paragraph{Protocol.} The unit of replication is the video. Each test video is
monitored alone by a freshly started detector, and no frame is reused between
calibration and monitoring. Calibration uses training videos only. For Ped2
and Avenue it uses all training videos of the same camera. For Mivia it uses
the training videos of the other folds, pooled across cameras.

Operating points are selected on training videos only, before any test video
is read. SENTINEL-E's betting lag $L^\ast(\alpha)$ is the smallest lag on the
grid $\{1,2,3,5,7,10,15,20,30,50\}$ whose leave-one-video-out (LOVO)
per-video false-alarm proportion is at most $\alpha$. Ped2 and Avenue use
LOVO over their 16 training videos. Mivia uses 10 folds by video index, each
holding out 20 training videos. Each classical baseline is thinned at the same
$L^\ast(\alpha)$. Its knob is the first admissible knob in sweep order, from the
most aggressive. A knob is admissible when its LOVO proportion is at most
$\alpha$. If no knob is admissible, the most conservative knob is used. This
happens for \RThreeBaselineNoKnob{} of the \RThreeBaselineConfigs{} baseline
configurations at $\alpha=0.1$, and those results describe a conservative
setting, not a tuned one.

This replaces an earlier evaluation in which the betting lag was tuned on test
streams and normal frames were reused across spliced streams.

A null trial is a video with no anomaly, or the pre-onset part of an anomalous
video that lasts at least 2\,s. On Mivia the null trials are the 50 negative
test videos. A detection is an alarm inside an episode window, extended by a
one-second grace period. The detector restarts after each alarm, so a video
can contribute more than one alarm to the rate per hour. It counts once in the
per-video proportion. All proportions use exact Clopper--Pearson intervals.
False alarms per hour use an exact Poisson interval. The effective sample size
is stated with each result: the number of null videos, and the number of
episodes. Avenue has 74 episodes in 21 videos, so its detection intervals treat
episodes as independent and are optimistic.

\paragraph{Results.} \Cref{tab:r3_selection} gives the selected lags, and
\cref{tab:r3_test} the test results at $\alpha=0.1$. \Cref{fig:r3_per_video}
plots each method's test false alarms per hour against its detection rate.

\inputtable{tabr3_selection}

\inputtable{tabr3_test}

\begin{figure}[t]
  \centering
  \resultfig{fig_r3_per_video}{}
  \caption{Test results at $\alpha=0.1$, one panel per dataset and backbone.
  Each point is one method at its training-selected knob and $L^\ast(0.1)$.
  The horizontal axis is test false alarms per hour on a log scale, with zero
  drawn at $0.1$ per hour. The vertical bars are 95\% Clopper--Pearson
  intervals for the episode detection rate.}
  \label{fig:r3_per_video}
\end{figure}

The first finding concerns validity. On Ped2 and Avenue the false-alarm counts at
$\alpha=0.1$ are small. The null footage is short: 12.6\,s for Ped2 (three pre-onset
segments), 6.9\,min for Avenue (16 videos) and 26.9\,min for Mivia (50 negative
videos). At this length any outcome is consistent with the nominal level. Ped2 flow has \RThreePedTwoFlowFAk{}/\RThreePedTwoFlowFAn{} false
alarms, Ped2 PatchCore has \RThreePedTwoPCFAk{}/\RThreePedTwoPCFAn{}, Avenue
flow has \RThreeAvenueFlowFAk{}/\RThreeAvenueFlowFAn{} and Avenue PatchCore has
\RThreeAvenuePCFAk{}/\RThreeAvenuePCFAn{}. The intervals are wide. Ped2 flow at
$\alpha=0.2$ has 3 of 3 false alarms, with selected lag 1 (Online Resource~1, Table~S1),
so the Ped2 result depends on the level chosen. Mivia flow has
\RThreeMiviaFlowFAk{}/\RThreeMiviaFlowFAn{} false alarms, inside the level.
Mivia PatchCore has \RThreeMiviaPCFAk{}/\RThreeMiviaPCFAn{}, a proportion of
\RThreeMiviaPCFAp{} with interval $[\RThreeMiviaPCFALo{},\RThreeMiviaPCFAHi{}]$.
This exceeds $0.1$ at the point estimate, and the interval contains $0.1$. The
one-sided exact binomial probability of at least nine false alarms in 50 videos at
$p=0.1$ is $0.058$, so the excess is not significant at the 5\% level. The count also
moves with $\alpha$: it is 3 of 50 at $\alpha=0.05$ and 12 of 50 at $\alpha=0.2$
(full sweep in Online Resource~1, Sect.~S1). The excess has several candidate causes,
and 50 null videos cannot separate them. First, Mivia calibrates across cameras, but
the guarantee is per camera, and every test camera is unseen, so the per-camera
premise is not tested here. Second, the p-values are serially dependent (the
paragraph below). Third, the lag and the alarm boundary are chosen on training videos.
The selected lag is the smallest one whose LOVO proportion is at most $\alpha$. For
Mivia PatchCore, $L^\ast=7$ has 20 of 200 $=0.100$, exactly at $\alpha$, and selection at
a boundary can pick a lag whose true rate lies above $\alpha$.

Per-video rates understate what an operator sees, because the detector restarts after
each alarm. Per hour of null footage, SENTINEL-E raises 284.8 false alarms on Ped2
flow, from one alarm in 12.6\,s. On Avenue PatchCore it raises 86.6 per hour, and all
ten of its alarms come from one null video in 6.9\,min. On Mivia PatchCore it raises
100.5 per hour, with 45 alarms across nine videos in 26.9\,min. A per-video
proportion counts a video once, however many alarms it carries. The hourly figures
rest on very short null footage and are illustrations, not operating rates.

The second finding concerns power, which depends strongly on the backbone.
Ped2 flow detects \RThreePedTwoFlowDetk{} of \RThreePedTwoFlowDetn{} episodes
and Ped2 PatchCore detects \RThreePedTwoPCDetk{} of \RThreePedTwoPCDetn{}.
Avenue detects only \RThreeAvenueFlowDetk{} of \RThreeAvenueFlowDetn{} with
flow and \RThreeAvenuePCDetk{} of \RThreeAvenuePCDetn{} with PatchCore. Mivia
flow detects none of its \RThreeMiviaFlowDetn{} positive videos. Mivia
PatchCore detects \RThreeMiviaPCDetk{} of \RThreeMiviaPCDetn{}. Its median
delay is \RThreeMiviaPCDelayMed{}\,s after onset.

The third finding concerns the baselines, and it does not favour SENTINEL-E. On
Mivia PatchCore every baseline also exceeds $0.1$ on test, with 9 to 18 of 50
negative videos raising an alarm. On Avenue PatchCore, three of the six baselines
(CUSUM, Shiryaev--Roberts and the parametric e-detector) had no admissible knob at
$\alpha=0.1$ and used the most conservative setting. Across all 36 baseline
configurations at this level, 15 had none. The fixed threshold and E-SHIFT detect
more episodes than SENTINEL-E on Avenue PatchCore (15 and 29 of 74, against 9), with
2 and 4 false alarms in 16 videos. On Mivia PatchCore, CUSUM matches the 9 false
alarms in 50 videos, detects more episodes (11 of 50 against 10) and has a lower rate
per hour (53.6 against 100.5). Its knob is the most conservative in its sweep, so
this is not a tuned baseline. On Ped2 PatchCore the fixed threshold detects 8 of 12
episodes against 7 of 12, with no false alarm in three null videos, which is too
little evidence to rule out an excess. SENTINEL-E therefore shows no demonstrated
advantage over these baselines on real video. Its claimed advantage is the
horizon-free guarantee, which the short real videos cannot exercise.

\paragraph{Dependence.} The guarantee rests on the thinned p-values, so
\cref{tab:r3_dependence} measures their dependence. Two quantities are
reported. The pooled autocorrelation of the concatenated training scores
exceeds $0.05$ at every lag up to 200. That length reflects differences
between the mean scores of videos, not the time scale of the process, and
it is not used. The within-video autocorrelation is the relevant quantity. At
lag one it averages \RThreeWithinLagOneMin{} to \RThreeWithinLagOneMax{}. At
lag 25 it averages between essentially zero and \RThreeWithinLagTwentyFiveMax{}.
Lag 25 is 1\,s at 25\,fps and 2.5\,s at 10\,fps.

The thinned p-values remain positively correlated. At the selected lags the
lag-one correlation ranges from \RThreeDepCorrMin{} to \RThreeDepCorrMax{}. The
empirical CDF shows no excess of small values. The largest positive deviation
is \RThreeDepSupMax{} at the selected lags, and \RThreeDepSupOneMax{} at lag one.
The conditional-mean diagnostic contradicts conditional super-uniformity. If $p_t$
is conditionally super-uniform given the past, its mean given $p_{t-1}\le0.1$ is at
least $0.5$. The observed means lie between \RThreeDepCondMin{} and
\RThreeDepCondMax{}, on \RThreeDepCountMin{} to \RThreeDepCountMax{} pairs. A small
p-value is followed by another small one. The hypothesis of \cref{thm:episodic}
therefore fails on these thinned streams, and real-video validity can only be
established empirically. The observable consequence is the per-video false-alarm
rate. On Ped2 and Avenue it is consistent with $0.1$, but the null footage is short.
On Mivia PatchCore it exceeds $0.1$.

\inputtable{tabr3_dependence}

\paragraph{Action or residue.} The Mivia labels mark only the onset. They do
not say whether the detector responds to the dumping action or to the waste
left behind. Flow measures motion. Flow scores are higher on negative videos
than before onset in positive videos, so flow can at best respond to the
action, and the motion of other activity in the scene confuses it. Appearance
can in principle see deposited waste. Under PatchCore, static dumping is
detected in \RThreeMiviaPCStaticDetk{} of \RThreeMiviaPCStaticDetn{} videos and
dynamic dumping in \RThreeMiviaPCDynDetk{} of \RThreeMiviaPCDynDetn{}. These
counts do not differ at this sample size. The per-type split for all methods
is in Online Resource 1, Sect.~S1. A residue that persists after the action is
the permanent-change case, in which the episode never ends, $\eta=0$. The
episodic mixture contains that case. A dedicated residue model would change
the background model itself, and we leave it to future work.

\paragraph{Scope.} The real test videos last seconds to about a minute. The real data therefore bear only weakly on per-video false-alarm control, because
the null footage is short. The long-horizon
claim rests on the simulated streams in \cref{sec:validity}, which run for up
to a week of simulated time.

\begin{figure}[t]
  \centering
  \resultfig{fig_r0_qualitative}{}
  \caption{Real frames, real scores. Ped2 and Avenue rows: one normal and one
  anomalous frame, with the Farneback flow-magnitude field and its score, at
  the $64\times64$ resolution the pipeline uses. Mivia rows: the static and
  dynamic test videos, each with one pre-onset frame, one post-onset frame and
  the PatchCore score over time. Each video is the first positive test video of
  its type, in video-id order, with at least 3\,s before onset. It was not
  chosen on its score. The dashed line marks onset, and the coloured points mark
  the two frames. The frame panels are shown at display resolution, from the
  original releases. The Mivia release is not redistributed, so only downscaled
  stills are stored. The flow and score panels are computed from the pipeline
  inputs.}
  \label{fig:qualitative}
\end{figure}

The qualitative panels show how small the per-frame signal is. Ped2 moves from
\QualExamplePedTwoNormal{} to \QualExamplePedTwoAnomalous{}, and Avenue moves
from \QualExampleAvenueNormal{} to \QualExampleAvenueAnomalous{}. These are
the values the pipeline consumes, and they explain why the method accumulates
evidence over time.


\subsection{A single stream, worked through}
\label{sec:worked-example}

Aggregates over hundreds of streams are hard to read without one stream in mind.
\Cref{fig:worked-example} shows a \WorkedT{}-second stream with
\WorkedNEpisodes{} episodes. It is the first seed tried at this configuration,
not one chosen for a fast detection.

\begin{figure}[t]
  \centering
  \resultfig{fig6_worked_example}{}
  \caption{One simulated \WorkedT{}-second stream with \WorkedNEpisodes{}
  episodes (shaded). \textbf{(a)} The raw backbone score. \textbf{(b)} The
  conformal p-value at each betting instant (log scale). Score spikes do not map
  onto uniformly small p-values, because the conformal layer absorbs much of the
  diurnal and weather variation. \textbf{(c)} The episodic log-wealth. The
  dash-dot line marks the labelled onset, and the red point marks the alarm at
  $\log(1/\alpha)$, \WorkedDelaySec{}\,s after onset.}
  \label{fig:worked-example}
\end{figure}

The wealth stays close to its starting value for most of the stream. Under the
null it is a supermartingale, so it drifts but has no reason to grow. Once the
first episode's evidence accumulates, the wealth climbs to the alarm level
$1/\alpha=100$, \WorkedDelaySec{}\,s after onset. The later episodes do not
move the trace, because this is a single-alarm illustration. The restarted,
multi-alarm behaviour is measured in the sections below.
\texttt{experiments/exp11\_worked\_example.py} regenerates the figure for
another seed.

\subsection{False-alarm control is time-uniform}
\label{sec:validity}

\begin{figure}[t]
  \centering
  \resultfig{fig1_validity}{}
  \caption{Time-uniform false-alarm control, \ExpOneReps{} null streams per point
  in (a) and \ExpOneHorizonReps{} per point in (b). Open downward triangles mark
  points with no observed false alarm, drawn at the 95\% Wilson upper limit.
  \textbf{(a)} Realised against nominal level. SENTINEL-E stays below the
  diagonal. The per-frame p-value threshold, fed identical conformal p-values,
  sits at the top of the plot. The middle curve is the same e-detector fed exactly
  uniform p-values. \textbf{(b)} Against the horizon at $\alpha=0.01$. The
  per-frame rules approach one, and SENTINEL-E stays flat.}
  \label{fig:validity}
\end{figure}

\inputtable{tab1_validity}

At every level in \cref{tab:validity}, the upper 95\% limit on SENTINEL-E's
realised false-alarm probability is at most \ValidityWorstCiHi{}. At the
tightest level, $\alpha=\ValidityTightestAlpha$, the rate is $0.000$ with upper
limit \ValidityTightestCiHi{}, which is below nominal. The per-frame threshold
on the same p-values reaches \ValidityPfaPvalue{} at $\alpha=0.01$. Its mean time
to the first false alarm is \ValidityPvalueArl{} frames, or
\ValidityPvalueArlMin{}\,min of video. Across a hundred cameras, that is a false
dispatch every few seconds.

\Cref{tab:horizon_week} and \cref{fig:horizon-week} extend the horizon to a week.
They use a simulated diurnal cycle of a full 24\,h of $2{,}160{,}000$ frames, and calibration
covers one full day. SENTINEL-E runs on \HorizonWeekReps{} streams per horizon,
and the two per-frame rules on \HorizonWeekBaselineReps{}, up to
\HorizonWeekMaxHours{}\,h or \HorizonWeekMaxMFrames{} million frames. SENTINEL-E
raised no false alarm in any stream at any horizon. Its 95\% upper limit on the
false-alarm probability is \HorizonWeekSentinelCiHi{}, which does not by itself
resolve $\alpha=0.01$. The p-value threshold
reaches \HorizonWeekPvalueOneHour{} within the first hour. The fixed threshold
reaches \HorizonWeekFixedOneHour{} at one hour and \HorizonWeekFixedSixHour{}
from six hours.

\begin{figure}[t]
  \centering
  \resultfig{fig_horizon_week}{}
  \caption{False-alarm probability against horizon over week-long null streams at
  $\alpha=0.01$, with a 24\,h diurnal cycle. Dashed line: nominal level.}
  \label{fig:horizon-week}
\end{figure}

\inputtable{tab_horizon_week}

SENTINEL-E is conservative, and the source of the conservatism is specific.
Feeding the same e-detector exactly uniform p-values gives a realised rate of
\ValidityPfaOracleLoose{} at nominal \ValidityAlphaLoose. Most of the gap is
therefore slack in Ville's inequality and in the prior, not in the conformal
layer. A practitioner who needs a rate close to nominal should calibrate the
threshold empirically. The theory provides something a fixed threshold cannot:
the bound does not degrade with the horizon.

\subsection{Episodes versus changepoints}
\label{sec:episodes}

Two mixtures are compared on the same seeds. The episodic arm mixes over six
values of $\kappa$ and four of $\eta$, giving \EpiGridPoints{} grid points. The
changepoint arm mixes over \EpiGridPoints{} values of $\kappa$ with $\eta=0$.
The two arms therefore use grids of equal size, 24 points: 6 $\kappa$ by 4 $\eta$
for the episodic arm and 24 $\kappa$ values for the changepoint arm. The grids are
not shared. The changepoint arm covers $\kappa$ more finely, which biases the
comparison against the episodic arm. Delay differences are reported as paired bootstrap intervals, and miss
differences use an exact McNemar test on discordant pairs. Unpaired comparisons
on different grids are in Online Resource~1, Sect.~S2. They are descriptive only.

\inputtable{tab13_paired_episodes}

The sign of the difference depends on the number of episodes. At one episode,
the episodic mixture is behind. The censored-delay difference is
\EpiPairedDiffOne\,s (95\% CI \EpiPairedCIOne{}, bootstrap $p<0.001$),
and the changepoint arm detects \EpiPairedDiscordantOne{} events the episodic arm
misses, with none in reverse (McNemar $p=\EpiPairedPMcnemarOne$). This is the
cost of the extra axis: prior mass goes to episode-end hypotheses that one
episode does not repay. At two and four episodes, the episodic mixture is ahead, by
\EpiPairedDiffTwo\,s and \EpiPairedDiffMid\,s. The discordant misses all favour
it, \EpiPairedDiscordantTwo{} and \EpiPairedDiscordantMid{} respectively, with
McNemar $p=\EpiPairedPMcnemarTwo$ and $p=\EpiPairedPMcnemarMid$. At eight episodes
the delay advantage is smaller, \EpiPairedDiffMany\,s (bootstrap
$p=\EpiPairedPbootMany$), and the miss difference is not significant
($p=\EpiPairedPMcnemarMany$).

The configuration matters too. \Cref{tab:corner_sweep} and \cref{fig:corner}
repeat the paired comparison at $\alpha=0.01$ across the event shapes the
simulator exposes, with \CornerReps{} seeded pairs each. A delay difference below 1\,s in absolute
value counts as equivalent. This is the margin used in \cref{tab:corner_sweep}.

\begin{figure}[t]
  \centering
  \resultfig{fig_corner_sweep}{}
  \caption{Censored-delay difference (changepoint minus episodic) across simulator
  configurations, with bootstrap intervals. Positive values favour the episodic
  mixture.}
  \label{fig:corner}
\end{figure}

\inputtable{tab_corner_sweep}

Long, intermittent episodes favour the episodic mixture. With 1600-frame episodes
at $\pi=0.2$ and a gap of $12{,}000$ frames, the episodic arm is ahead by
\CornerLongGapDiff\,s (CI \CornerLongGapCi), with \CornerLongGapEpiOnly{}
discordant misses and none in reverse. With a gap of $1500$ frames the delay
advantage is \CornerShortGapDiff\,s (CI \CornerShortGapCi), but three discordant
misses give McNemar $p=\CornerShortGapP$, so the miss difference is not resolved.
Dense episodes ($\pi=0.7$) and a permanent change have delays of 9 to 12\,s. The
changepoint mixture is faster by \CornerDenseAbsDiff\,s in both cases. The 95\%
interval of that difference lies within 0.2\,s of zero, so both rows are equivalent
under the 1\,s margin and the difference is negligible. A single episode of 800
frames favours the changepoint mixture by \CornerSingleAbsDiff\,s (CI
\CornerSingleCi). Short episodes of 200 frames are missed by both mixtures at
$\alpha=0.01$, a miss rate of \CornerFloorMiss{}, which marks the detection floor
for that event length.

\label{sec:per-episode}
Stream-level scores can hide episode-level misses. \Cref{tab:per_episode} scores
each episode's own window, using \ExpTenReps{} streams per row. At four episodes the
stream miss rate is \PerEpiMidEpiStreamMiss{}, but the per-episode miss rate is
\PerEpiMidEpiMiss{}. At eight episodes, the stream miss rate is
\PerEpiManyEpiStreamMiss{} and the per-episode rate is \PerEpiManyEpiMiss{}.
The two mixtures have indistinguishable per-episode miss rates at every count:
\PerEpiMidEpiMiss{} against \PerEpiMidCpMiss{} at four episodes. The difference
lies in how quickly evidence that does not resolve inside one episode
accumulates across the stream. The paired tests above measure that, not
independent catches of each episode.

\inputtable{tab14_per_episode}

\subsection{Detection delay at a controlled false-alarm rate}
\label{sec:delay}

\begin{figure}[t]
  \centering
  \resultfig{fig2_delay_far}{}
  \caption{Detection delay against realised false-alarm rate, \ExpTwoReps{} null
  and \ExpTwoReps{} post-change streams per operating point. Each method is swept
  over its own knob and plotted where it lands. Lower and further left is better.}
  \label{fig:delayfar}
\end{figure}

\inputtable{tab3_delay_far}

\Cref{tab:delay_far} gives the delay of each method at a controlled false-alarm
rate. Baselines are shown at the tightest realised false-alarm rate that still leaves
some events detected. The miss-rate column shows what that costs. The two exact e-processes are compared at a common nominal
level, because their nominal level is the knob that an operator sets.
Each $\pm$ is 1.96 standard errors. At $\alpha=0.2$ the two tie: SENTINEL-E has a
censored delay of \EqSEFifthDelaySec\,s ($\pm\EqSEFifthCiSec$) and the changepoint
mixture \EqCpFifthDelaySec\,s. At $\alpha=0.1$ the changepoint mixture is faster,
at 138.1\,s against 146.4\,s for SENTINEL-E, with intervals of about $\pm52$\,s for
each. At $\alpha=0.05$ the figures are 181.4\,s for the changepoint mixture and
171.2\,s for SENTINEL-E. At $\alpha=0.01$, SENTINEL-E is faster, at
\EqSEHundredthDelaySec\,s against \EqCpHundredthDelaySec\,s, with miss rate
\EqSEHundredthMiss{} against \EqCpHundredthMiss. The ordering therefore changes
with $\alpha$, and the differences at $\alpha=0.2$ and $0.1$ are small against
their intervals. The nominal level
$\alpha=0.5$ is not a usable operating point. Its bound allows a false alarm on
half of all streams, so the low delay it reports (\EqSEHalfDelaySec\,s) says
little about a usable threshold.

The baselines behave differently. \DelayNumInadmissible{} of them never reach a
realised false-alarm probability of $0.05$ at any knob setting while still
detecting events (\DelayInadmissibleList). On streams with drift, heavy tails and
autocorrelation, their thresholds rest on assumptions that these streams violate, so they cannot be tuned into the operating region.

\subsection{Robustness to calibration/deployment mismatch}
\label{sec:shift}

\begin{figure}[t]
  \centering
  \resultfig{fig3_shift}{}
  \caption{Shift in an \emph{unbinned} covariate (scene activity). Mondrian calibration tracks the pooled baseline, because binning cannot help
  with a variable outside the taxonomy. Conformalising the context-normalised
  residual holds throughout.}
  \label{fig:shift}
\end{figure}

\inputtable{tab5_activity}

Two mismatches are routine between the week a calibration set is labelled and the
month a camera is monitored. An unbinned covariate can drift, and calibration can
cover only some of the conditions. \Cref{tab:activity} shows the first. At
$\times 3.0$ activity, pooled and Mondrian calibration both reach a false-alarm
probability of \ShiftPooledPfa{}, because activity is not one of the bins. The
context-normalised residual construction does far better, but at \ExpThreeReps{} streams per
row several shift levels cannot be resolved (\ShiftResidualUnresolvedRatios{}).

\inputtable{tab_shift_remedy}

\Cref{tab:shift_remedy} gives the sweep at 1,500 null streams and 200 post-change streams per level. The unweighted residual default is valid at $\times 1.0$ and $\times 1.5$, with a false-alarm probability of 0.0020 at $\times 1.5$. It is unresolved at $\times 2.0$ (0.0140, interval [0.0092, 0.0213]) and invalid at $\times 3.0$ (0.0620, interval [0.0509, 0.0754]). False alarms cluster immediately after the detector re-enters the calibrated context region, at frames near the support edge where the residual model extrapolates.

Likelihood-ratio weighting holds the level at every shift tested, but its power falls. Its miss rate is 0.095 at $\times 1.0$, against 0.025 unweighted, and rises to 0.540, 0.930 and 1.000 at $\times 1.5$, $\times 2.0$ and $\times 3.0$. The Kish effective-sample-size substitution has no calibration-conditional proof, so the weighted path is a heuristic. The pre-specified rule, reported in Online Resource~1, Table~S9, recommends weighting only if its extra miss rate at $\times 1.0$ is at most 0.05. The extra rate is 0.070, so the unweighted residual stays the default.

The held-out context gate, which would switch weighting on, is a classifier two-sample test with AUC at least 0.60. It fired on every null stream at every level, including the unshifted $\times 1.0$. Context tests respond to benign diurnal variation, so we use a stricter support region instead.

A stricter support region addresses the boundary effect. The residual calibrator abstains on a frame whose context falls outside a per-coordinate box, or whose Mahalanobis distance exceeds a slack multiple of the calibration 0.999 quantile. The default multiplier is 1.5. We set it to 0.5, so boundary frames abstain. An abstained frame receives a neutral bet, so the supermartingale argument is unchanged. The stricter region holds the level at every shift tested. At $\times 3.0$ the false-alarm probability falls from 0.0620 to 0.0007 (interval [0.0001, 0.0038]), and the abstention rate on null betting instants rises to 0.500. At $\times 1.0$ the miss rate equals the default, 0.025, and abstention is 0.010. Likelihood-ratio weighting is also valid, but its miss rate is 0.540 to 1.000 at $\times 1.5$ to $\times 3.0$, against 0.070 to 0.450 for the stricter region.

The slack value 0.5 was selected in a 300-stream pilot of seven variants (cool-down, box-margin and slack settings) whose seeds overlap the main run; only slack 0.5 and a cool-down at $k=50$ met the pre-set criterion. We confirmed slack 0.5 on 1,500 null streams and 200 post-change streams per level, with seeds disjoint from both the pilot and the main run. The evidence covers one shift family, a multiplier on scene activity, in one direction. The false-alarm probabilities are 0.0000, 0.0007, 0.0007 and 0.0013 at $\times 1.0$, $\times 1.5$, $\times 2.0$ and $\times 3.0$, and all are valid. The miss rates are 0.010, 0.040, 0.140 and 0.470.

The abstention rate is the operator-visible shift monitor. With slack 0.5 it is 0.010 at $\times 1.0$, 0.068 at $\times 1.5$ and 0.500 at $\times 3.0$. As a suggestion, not a validated rule, recalibrate when the abstention fraction over the monitoring window exceeds 0.05. That threshold separates the unshifted and $\times 1.5$ cases in this sweep, but no other rule or shift was tested. We recommend slack 0.5 for deployment. Without shift it matched the default's miss rate, and it was valid at every shift tested, whereas the default 1.5 is invalid at $\times 3.0$. Every other experiment in this paper used 1.5, and the released code keeps 1.5 as its default only so that those results reproduce exactly.

\Cref{tab:coverage} varies coverage instead. With calibration from one daytime
window, the residual construction keeps its level, but its miss rate reaches
\CoverDaytimeMiss{} and \CoverDaytimeAbstain\% of frames are abstained on. The
degradation appears in the abstention and miss columns, where an operator can see
it, rather than in the alarm rate. The calibration set must therefore span the
conditions to be monitored, and the abstention rate shows whether it does.

\inputtable{tab4_coverage}

\subsection{The camera network}
\label{sec:network}

\begin{figure}[t]
  \centering
  \resultfig{fig4_network}{}
  \caption{Fleets of \NetCameras{} cameras, with \NetFleets{} held-out fleets with
  events and \NetFleets{} all-null fleets. The learned spatial prior does not
  shorten delay. e-BH holds the false-discovery rate under the dependence the
  graph induces, whichever controller feeds it.}
  \label{fig:network}
\end{figure}

\inputtable{tab6_network}

The fleet has \NetCameras{} cameras. Each evaluation fleet covers \NetEvalMinutes{}
minutes ($60{,}000$ frames) with three episodes of 800 frames per camera, at
$\pi=0.35$ and a mean gap of $6000$ frames. An episode propagates to a neighbour
with probability \NetPropProb{} and a mean lag of \NetPropLagSec{}\,s. The spatial
prior is learned from \NetTrainFleets{} training fleets of $30{,}000$ frames.
Evaluation uses \NetFleets{} fleets with events and \NetFleets{} all-null fleets,
with $\alpha=0.01$ per camera and e-BH at level \NetFdrLevel{}. This configuration
is not the one in \cref{sec:delay}, which has four episodes in 60 minutes. The
independent detectors miss \NetNoGraphMiss{} of events here, not the
\DelaySentinelMiss{} of \cref{sec:delay}.

\paragraph{The guarantee survives learned control.}
This is what \cref{thm:predictable} says, and the simulation agrees. On all-null
fleets, every controller (none, hand-designed and learned) gives an alarm rate of
at most \NetGnnNullPfa{} and an e-BH false-discovery rate of \NetGnnFdr{}, against
a target of \NetFdrLevel{}. The learned controller reduces power, and it still
leaves the false-alarm guarantee intact, because it can act only through a
predictable modulation.

\paragraph{The learned spatial prior lengthens delay.}
The hand-designed prior changes delay only slightly, and the learned prior lengthens it. Conditional delay is \NetHeuristicDelaySec\,s for the hand-designed prior and \NetGnnDelaySec\,s for the learned one, against \NetNoGraphDelaySec\,s with no graph. Censored delay is \NetHeuristicCensDelaySec\,s
and \NetGnnCensDelaySec\,s, against \NetNoGraphCensDelaySec\,s. Power is
\NetHeuristicPower{} for the hand-designed prior, \NetNoGraphPower{} with no graph,
and \NetGnnPower{} for the learned controller. Relative to no graph, the learned controller has lower power and longer delay.

Two controller variants were compared. Initialising the heads to reproduce
the ungoverned detector is now the default, and it is what \cref{tab:network}
reports. Crediting log-wealth over a window of \NetCreditWindowSec\,s from each camera's onset, in place of terminal log-wealth, lowered power from
\NetGnnPower{} to \NetGnnWindowedPower{}. Both configurations are reported, and
the windowed objective stays in the library behind a flag. The controller has
little signal to learn from. A denser fleet, stronger propagation or a shared
access road might change this, but these were not tested.

\paragraph{Fleet-level FDR control does earn its place.}
e-BH recovers the affected cameras at power \NetNoGraphPower{}, with a realised
false-discovery rate of \NetNoGraphFdr{} against a target of \NetFdrLevel{}. The
dependence comes from shared weather, shared lighting and correlated episodes.
In these fleets the power comes from the e-BH combination, and the hand-designed prior changes it by half a percentage point. This part of the fleet layer is the part worth keeping.

\subsection{Computational cost}
\label{sec:compute}

\inputtable{tab_cost_breakdown}

The conformal p-value and the e-detector step run only at betting instants,
one per decorrelation lag. With $L=\CostLag{}$, the batch path costs
$(0.34+13.69)/20 = 0.702$\,\textmu s per frame predicted, with the measured cost
given below. The streaming path, which a deployed camera runs, costs
$(0.34+39.28)/20 = 1.98$\,\textmu s per frame predicted. The guarantee layer therefore
costs under two microseconds per frame on either path. The lag $L=20$ is the
decorrelation lag estimated on the 200{,}000-frame benchmark stream. The simulation
experiments use a lag of 18, and the lags selected on real video range from 1 to
30. Because the batch cost scales as $1/L$, a lag of 1 gives about 14\,\textmu s per
frame.

The measured end-to-end cost is $0.610$\,\textmu s per frame, or 0.122\,s over a
200,000-frame stream. It includes the $O(T)$ expansion of the bets onto the
frame grid, which costs 0.017\,\textmu s per frame. The prediction is about 15\%
above the measurement, because the per-call component timings carry overhead
that one long run amortises. Against the 36.84\,ms frozen backbone, the measured
cost is 0.0017\% per frame, and the predicted cost is 0.0019\%.

The figures exclude context-feature extraction, video decoding, the backbone
itself and peak memory. The backbone is a stand-in for an edge detector, with
\CostBackboneParams{}\,M parameters at $320\times320$, not a published model.
Fitting a camera's calibration takes \CostFitSec{}\,s, once.

\subsection{Checking the assumptions}
\label{sec:diagnostics}

\inputtable{tab9_diagnostics}
\inputtable{tab10_pcoverage}

The guarantee rests on two premises. The first is conditional super-uniformity of
the betting increments. The second is marginal calibration of the p-values, which is checked in simulation
below. The first is testable, and the conditional-mean test on real video departs from it (\cref{sec:real-data}). The operational consequence of both is the
per-video false-alarm rate measured on real video in \cref{sec:real-data}.

\Cref{tab:diagnostics} shows the conditional-independence diagnostics on the
calibration residuals of \DiagStreams{} simulated cameras. Betting on every frame leaves a
lag-one autocorrelation of \DiagAcfRaw{} and a median Ljung--Box $p$ of
\DiagLjungRaw. Betting once per estimated lag, at a mean lag of \DiagLagMean{}
frames, gives \DiagAcfThin{} and a median $p$ of \DiagLjungThin{}. Thinning reduces
the violation but does not remove it. Only \DiagLjungFrac{} of cameras pass the
Ljung--Box test at 5\%. The diagnostic therefore detects residual dependence after thinning. What the experiments support is the empirical false-alarm control of
\cref{sec:validity} and \cref{sec:real-data}, not the literal hypothesis of
\cref{thm:episodic}.

\Cref{tab:pcoverage} checks the marginal calibration of the deployment p-values.
Over \DiagPN{} betting instants pooled from labelled null streams, the p-values
are super-uniform at every level tested. The smallest observed p-value is
\DiagPmin{}. At a calibration size of \CalN{}, the exact Beta level is
\BetaPmin{}, against \DkwPmin{} for the DKW inflation, a factor of
\BetaOverDkwSpeedup{}. The check uses labelled null streams, which an operator
does not have. The thinned calibration set is the quantity an operator can compute.

\subsection{Ablation}
\label{sec:ablation}

\inputtable{tab_ablation_highrep}

At 1,500 null streams per row, all seven load-bearing rows resolve. Four are
valid and three are invalid at $\alpha=0.01$.

The valid rows are as follows. The full SENTINEL-E has realised PFA
$0.0000$ [0.0000, 0.0026] and a censored delay of 244.5\,s. The changepoint
mixture ($\eta=0$) is also valid, with PFA $0.0000$ [0.0000, 0.0026]. It misses
0.055 of events against 0.015 for the episodic mixture, and its censored delay
is 311.3\,s. The DKW inflation is valid with PFA $0.0000$ [0.0000, 0.0026]. It
misses 0.675 of events, with a censored delay of 2025.8\,s, against 244.5\,s for
the exact Beta levels. Both rows are valid, so the gap is a difference in power
at the same level, not a trade-off against validity.

Unthinned calibration is valid too, with PFA $0.0020$ [0.0007, 0.0059]. Its
censored delay is 68.4\,s, against 244.5\,s for the thinned calibration set.
\Cref{thm:beta} assumes i.i.d.\ calibration draws, and thinning by the
decorrelation lag is what supplies that premise. The experiment does not show
that thinning is necessary for validity.

The invalid rows are as follows. Betting on every frame gives PFA $1.0000$
[0.9974, 1.0000], so every null stream alarms, with a mean time to false alarm
of 337 frames. Pooled conformal calibration gives PFA $0.9887$ [0.9819, 0.9929].
Omitting the calibration-conditional correction ($\delta=0$) gives PFA
$0.0387$ [0.0300, 0.0497], nearly four times nominal. The correction is therefore required, and pooled calibration is not a substitute for it.

The other twelve variants of the 200-stream table are descriptive. They are
listed in Online Resource~1, Table~S8. At 200 streams, 14 of the 19 variants are
unresolved, so that table gives point estimates, not verdicts. Two of the twelve
resolve as invalid at that count: bet once per five frames (PFA 0.935) and
Mondrian conformal (0.025 [0.011, 0.057]). The fixed-start variant, which has no
changepoint prior, misses 1.00 of events. A wealth process started at frame one
has spent too much when a late change arrives. The mixture over onsets in
\cref{eq:episodic} avoids that outcome.

\section{Limitations and future work}
\label{sec:limitations}

\paragraph{Public benchmarks are short.}
The real test videos last seconds to minutes. The null footage is 12.6\,s for Ped2, 6.9\,min for Avenue and 26.9\,min for the Mivia-IWDD negative videos. Mivia-IWDD provides no per-camera calibration footage, so every test camera is unseen and calibration pools cameras. Pooling is one of three candidate explanations for the Mivia-IWDD PatchCore proportion of 9 of 50 normal videos (\cref{sec:real-data}), and per-camera calibration footage would test it directly. Ped2 and Avenue are pedestrian-walkway corpora rather than dumping footage, which is why Mivia-IWDD was added. Long-horizon and per-camera validation therefore await operational multi-week data. The week-long simulation rests on diurnal and weather assumptions that such data would test.

\paragraph{Backbones are frozen and training-free.}
The backbone study uses frozen, training-free backbones. The available compute (a single cloud GPU) and the need for a training-only selection protocol for per-dataset trained detectors both favoured this choice. Detection power depends on the backbone. Optical flow detects 11 of 12 Ped2 episodes, while PatchCore detects 9 of 74 Avenue episodes and 10 of 50 Mivia-IWDD dumping videos, where flow detects 3 and none. Training state-of-the-art, per-dataset video anomaly detectors under the same leak-free, training-only protocol is future work and the natural next step for power. The Mivia-IWDD labels mark only the onset, so the detections cannot yet separate the dumping action from the waste left behind. A residue-specific model is also future work.

\paragraph{The theory assumes conditionally super-uniform p-values.}
\Cref{thm:episodic} requires the betting p-values to be conditionally super-uniform given the past. On real video the thinned p-values retain serial dependence. The lag-one correlation lies between 0.758 and 0.985 at the selected lags, and after a p-value below 0.1 the next p-value has mean 0.015 to 0.203, against 0.5 under the premise (\cref{tab:r3_dependence}). On real video the time-uniform bound is therefore used as an operating rule, and its realised per-video false-alarm rates are verified empirically. Extending the theory to dependent p-values, for example through block betting, is future work. In simulation, thinning by the decorrelation lag reduces the dependence but does not remove it, and the Ljung--Box test rejects independence for most cameras \citep{ljung1978}.

\paragraph{The shift study covers one family.}
The stricter support region, with slack 0.5, holds its level at every shift tested, and it was confirmed on fresh seeds. The evidence covers one activity-shift family, in one direction. For other shifts the abstention rate is the monitor we propose. The likelihood-ratio weighted path has no calibration-conditional proof, because its Kish effective-sample-size substitution is a heuristic that \cref{thm:beta} does not cover. The recalibration threshold of 0.05 on the abstention rate is a suggestion, not a validated rule.

\paragraph{The episodic advantage is configuration-specific.}
Episodes help when they are long, intermittent and recurring (\cref{tab:corner_sweep}). With 1600-frame episodes, $\pi=0.2$ and a 12{,}000-frame gap, the episodic mixture is ahead by 141.1\,s (95\% CI +79.6 to +210.2\,s). A single 800-frame episode favours the changepoint mixture by 55.7\,s (95\% CI $-111.2$ to $-14.0$\,s), dense and permanent changes are equivalent under the 1\,s margin of \cref{sec:episodes}, and both mixtures miss 200-frame episodes at $\alpha=0.01$. The delay comparison and the ablation use one configuration, four episodes in 60 minutes (\cref{sec:delay}). The guarantee does not depend on the episode model, but power does. Heavy-tailed episode durations, which dumping may produce, are untested.

\paragraph{Deployment considerations.}
The time-uniform guarantee is the horizon-free counterpart of the in-control average run length, and an alarm budget per monitoring period is its operational counterpart. Control charts also need a recalibration cadence, which statistical process control treats as a design choice \citep{montgomery2020introduction}, and an ROC operating point \citep{fawcett2006introduction} describes the trade-off that control rooms already use. Hourly alarm rates on real video are large, at 100.5 per hour of null footage for PatchCore on Mivia-IWDD, so an alarm budget should be set against them as well as against per-video rates. The choice of $\alpha$, against the costs of a false dispatch and of a missed event, is left to each operator. Every alarm should go to a named human reviewer before any action. A deployment also needs a retention policy and a route by which residents can contest a flagged event, because public-space surveillance raises accountability questions that the guarantee does not answer.

\section{Conclusion}
\label{sec:conclusion}

Continuous monitoring needs its own evaluation. A per-frame rule fed valid p-values raises a false alarm on \ValidityPfaPvaluePct\% of 40-minute null streams, so a clip-level level does not carry over to a stream. This paper makes three contributions.

The theory is an episodic e-process. It mixes over episodes that start, end and recur. It is a non-negative supermartingale when the p-values are conditionally super-uniform, so Ville's inequality gives time-uniform false-alarm control, and it contains the changepoint mixture as the case in which episodes never end. On real video the conditional premise is not met after thinning, so validity there is checked only empirically (\cref{sec:limitations}).

The algorithm is a two-state forward recursion. It runs in constant time per frame, with the p-value and the e-detector step computed once per betting lag. The guarantee layer costs \CostSentinelUs\,\textmu s per frame in batch and about 2\,\textmu s per frame when streaming (1.98\,\textmu s predicted).

The evaluation reports realised rather than nominal false-alarm rates, censored delays and paired comparisons on grids of equal size.

The simulated results cover horizons of up to one week. No false alarm occurred in 200 week-long null streams (95\% upper limit 0.019, which does not by itself resolve $\alpha=0.01$), and per-frame thresholds reach 0.980 within an hour. The episodic mixture shortens delay for long, intermittent and recurring events. For a single change the changepoint mixture is faster, and for a permanent change the two are equivalent.

The real-video results are more limited in scale. On UCSD Ped2, CUHK Avenue and Mivia-IWDD, the per-video false-alarm proportions are consistent with the nominal level at $\alpha=0.1$ on all three corpora. The null footage is short (12.6\,s and 6.9\,min for Ped2 and Avenue). On Mivia-IWDD, PatchCore raises alarms on 9 of 50 normal videos, with a 95\% interval of [0.086, 0.314], and the supervised backbone raises alarms on 4 of 50. Alarm rates per hour of null footage are large: 100.5 per hour for PatchCore and 60.3 for the supervised backbone. Power depends on the backbone. Optical flow detects 11 of 12 Ped2 episodes and 3 of 74 Avenue episodes, PatchCore detects 7 of 12 and 9 of 74, and on Mivia-IWDD PatchCore detects 10 of 50 dumping videos and the supervised backbone 31 of 50. None of the backbones is state of the art. On Ped2 and Avenue, several classical baselines selected by the same training-only rule detect more episodes, generally with more false alarms. On Mivia-IWDD with the supervised backbone, SENTINEL-E ties with the p-value threshold for the most detections among methods that hold the nominal level. Its advantage is the horizon-free guarantee, which the short videos cannot test.

Three directions follow. The first is a per-dataset trained detector for Ped2 and Avenue under a training-only protocol, which would test power on those corpora without changing the theory. The second is per-camera calibration on Mivia-IWDD, which would test the guarantee under the condition it needs and show whether pooling accounts for the Mivia-IWDD proportion. The third is longer real streams, which would test the week-long simulation directly. Residue-specific change models would also separate the dumping action from the waste left behind. The argument and the algorithm do not depend on dumping, so the same optional-stopping problem applies to wildlife poaching, illegal fishing and infrastructure faults.


\section*{Declarations}

\paragraph{Competing interests.} The author declares no competing interests.
\paragraph{Data and code availability.} The real-video evaluation uses three public corpora: UCSD Ped2, CUHK Avenue, and Mivia-IWDD-500 (\url{https://github.com/MiviaLab/IllegalWasteDumping}). The backbone is a frozen, ImageNet-pretrained WideResNet-50-2 with torchvision weights. The corpora are distributed by their authors under their own terms, and no video is redistributed here, apart from a few illustrative still frames in the qualitative figure. Other corpora cited in Related Work were not used. All code, experiment scripts and generated results will be made public on acceptance. A private, non-identifying link is available to editors and reviewers on request during review.
\paragraph{Ethical approval.} Not applicable. No human-subjects data was collected for this work. The real video corpora (UCSD Ped2, CUHK Avenue and Mivia-IWDD-500) are publicly released benchmarks distributed by their original authors for research use. This paper is a secondary use of that public release, not new data collection, and involves no interaction with human subjects. \cref{sec:limitations} discusses the human-oversight, accountability and disclosure questions that automated surveillance of public space raises.
\paragraph{Funding and author contributions.} Provided on the separate title page accompanying this submission, per this venue's double-anonymous review policy.

\bibliography{refs}

\end{document}
